% !TeX root = ../../brightPhD.tex
\chapter{基于高阶有限元的多分辨率高精度拓扑优化}
\label{chap:high_fem_mtop}

%\iffalse

\section{引言}
\label{sec:ch4_intro}

在连续体结构拓扑优化中，获取边界清晰、物理响应准确且无数值伪影的优化构型一直是该领域的核心追求。在前一章中，本文系统探讨了高阶拉格朗日有限元在提升空间离散精度和物理场逼近能力方面的优势。然而，在传统的基于单元的单分辨率拓扑优化（Single-resolution Topology Optimization, STOP）框架下，分析网格与设计网格通常被刚性绑定，即一个分析单元对应一个设计变量。由此带来的直接后果是：即便采用高阶有限元改善了状态方程的离散精度，设计变量的几何分辨率仍然受限于分析单元的物理尺度。若进一步试图通过全局加密分析网格来提升拓扑细节表达能力，则往往会使系统自由度与计算规模迅速增长，从而形成“高精度分析”与“高分辨率设计”之间难以兼顾的突出矛盾。

为突破传统框架中分析分辨率与设计分辨率之间的刚性耦合，本章在 Nguyen 等人~\cite{nguyenComputationalParadigmMultiresolution2010a} 提出的 MTOP 基本范式的基础上，构建一种基于高阶有限元的多分辨率高精度计算框架。该框架通过引入相互独立且逻辑耦合的三层离散体系——位移有限元网格、密度积分网格与设计变量网格，实现了在相对较粗的高阶位移网格上，嵌套极高分辨率的材料分布表达。这一机制不仅充分发挥了高阶有限元在力学求解上的高精度与低自由度优势，同时赋予了优化算法捕捉复杂拓扑演化、精细特征分支以及平滑结构边界的能力。

在构建该框架后，首先通过基准算例验证了多分辨率策略在不同有限元阶次下的构型一致性与收敛鲁棒性。在此基础上，为抑制优化结果中的灰度过渡并增强构型的二值性，本章进一步探讨了不同密度映射与正则化机制（包括灵敏度过滤、密度过滤以及非线性投影算子）在多分辨率框架下的数学列式与作用机理。更为重要的是，针对拓扑优化中极具挑战性的局部应力约束问题，本章充分利用了高阶有限元对复杂应力梯度的高精度解析能力。结合多项式消失约束技术与 ALM，构建了面向大规模局部应力约束的 MTOP 求解算法，从而在高分辨率拓扑设计下实现了对局部强度的稳健控制。


\section{设计模型与分析模型}
\label{sec:design_analysis_model}

传统基于单元的变密度拓扑优化方法，通常将材料分布的参数化模型（设计模型）与偏微分方程的离散求解模型（分析模型）刚性绑定在同一离散网格上。这种处理方式在实现上虽然简洁直观，但掩盖了二者在数学功能与分辨率需求上的本质差异。事实上，从偏微分方程约束优化的底层逻辑来看，整个计算过程由两个本质不同但紧密耦合的子模型构成：设计模型负责生成满足正则化条件的可行密度场，决定拓扑表达的几何分辨率；分析模型则用于在给定密度场下高精度求解状态方程，决定力学响应的逼近精度。


\subsection{模型的定义与二分性}
\label{subsec:model_definition_binarity}

在连续层面，变密度拓扑优化的一般形式及其函数空间设定已在第~\ref{subsec:general_formulation} 节中给出。本文在本章中沿用相同的符号约定：以标量密度场 $\rho:\Omega\rightarrow[0,1]$ 描述结构中的材料分布，以位移场 $\boldsymbol{u}:\Omega\rightarrow\mathbb{R}^{d}$ 描述结构在外载与边界条件作用下的力学响应。

需要强调的是，从偏微分方程约束优化的严密视角来看，$\rho$ 与 $\boldsymbol{u}$ 并非地位对等的独立未知量，而是分别被定义为该优化系统中的控制变量与状态变量。二者通过材料插值模型与线弹性平衡方程形成明确的依赖映射：密度场 $\rho$ 决定了椭圆型偏微分算子的刚度系数，从而确定了状态方程的参数化形式；位移场 $\boldsymbol{u}$ 则是该边值问题在给定 $\rho$ 下的解，其隐式地刻画了该拓扑下的力学行为，并为结构响应的评估提供基础。

基于上述 “设计--分析” 的二分性认识，可将变密度拓扑优化的整个计算过程严格解耦为两个本质不同但逻辑关联的子模型：
\begin{itemize}
	\item \textbf{设计模型}：负责可行密度场的参数化表示。该模型通过插值、过滤与投影等映射机制，将设计变量转化为用于组装算子的物理密度场 $\rho$，实现对结构构型的几何描述。
	
	\item \textbf{分析模型}：负责在给定物理密度场的前提下，对线弹性状态方程进行空间离散与高精度求解。该模型通过评估位移场 $\boldsymbol{u}$ 及其派生量，实现对柔顺度、应力等目标与约束泛函的数值评估。
\end{itemize}

在离散层面，设计模型通常被参数化为一组有限维的设计变量向量 $\boldsymbol{d}$（如单元密度、节点密度或子单元密度等）。为严格区分“设计自由度”与“物理分析密度场”，可将物理密度的生成过程抽象为一条明确的数学映射链：首先由 $\boldsymbol{d}$ 经过空间插值构建连续的伪密度场近似 $\rho_h(\boldsymbol{x};\boldsymbol{d})$，随后通过过滤与投影等正则化算子，将其映射为具有明确物理意义的物理密度场 $\rho(\boldsymbol{x})$，进而参与有限元数值积分，组装形成全局刚度矩阵 $\boldsymbol{K}(\rho)$。相应地，分析模型在接收到 $\rho$ 后，通过求解离散的系统平衡方程：
\[
\boldsymbol{K}(\rho)\boldsymbol{U} = \boldsymbol{F},
\]
得到全局位移向量 $\boldsymbol{U}$，并据此评估目标泛函与约束条件，为后续灵敏度分析与优化变量更新提供基础。

从数值分析与工程设计的角度出发，在变密度方法中，设计模型与分析模型在分辨率需求上存在天然的差异：
\begin{enumerate}
	\item 为获得边界清晰、细节丰富且具有可控最小特征尺度的结构形态，设计模型倾向于采用更高的空间分辨率(即更密集的设计自由度)，以便刻画复杂的材料分布与细小构件。需要指出的是，尽管结构的物理最小特征尺度主要由过滤或投影等正则化算子的参数主导，但设计网格的空间分辨率直接决定了拓扑细节可表达的理论上限。
	
	\item 为了控制求解规模并避免状态方程的过度离散，分析模型通常倾向于在保证数值求解精度的前提下选取相对较粗的有限元网格，或是借助高阶有限元的强逼近能力在粗网格上进行近似，从而在较低的系统自由度下完成高精度的力学响应评估。
\end{enumerate}

综上所述，理想的拓扑优化框架应在空间离散层面实现设计与分析模型的相对独立：设计网格独立承载设计变量以生成连续密度场，分析网格则专注状态方程的高精度求解。这种数值层面的解耦不仅能充分释放高阶有限元的逼近优势，更为构建多分辨率计算框架奠定了基础。基于这一逻辑，本章不仅引入了相互独立的设计与分析网格，更进一步引入了高分辨率的密度积分网格，以确保刚度组装与局部约束评估的数值一致性。


\subsection{单分辨率框架的局限与精度错配}
\label{subsec:single_resolution_mismatch}

在传统 STOP 框架中，普遍采用离散的单元密度法：即为每个分析单元分配一个独立的设计变量。这种设计模型与分析模型在同一套网格上完全重合的策略，尽管在程序实现上简单直观，但却导致设计模型的几何分辨率被分析网格的物理尺寸刚性锁定，从而限制了二者在空间离散上独立调控的自由度。

这种强耦合不仅在概念上模糊了控制变量与状态变量的本质区别，更在数值层面上引发了显著的“精度错配”问题。当引入高阶有限元（阶次 $k \ge 2$）以期提升优化性能时，分析模型虽能在较粗网格上实现物理场的高精度计算，但受限于分片常数的密度离散假设，设计空间的拓扑表达能力并未随之提升。这意味着，高阶单元的高精度求解仅仅是对低分辨率阶梯状几何边界的冗余评估。这种精度的极度不对等不仅无益于刻画精细的拓扑特征，反而造成了计算资源在状态求解环节的严重浪费。

为直观揭示 STOP 框架中的精度错配现象，本文沿用第~\ref{subsec:benchmark_pathology_min_lengthscale} 节定义的对称二维 MBB 梁算例进行基准对比分析。在 STOP 中，设计变量与分析单元绑定（即设计变量自由度数 $\rho_{\text{dof}}=300$），在保持宏观分析网格（$30\times10$ 四边形网格）不变的前提下，若将位移有限元阶次从 $k=2$ 提升至 $k=4$，系统的位移自由度 $u_{\text{dof}}$ 将从 2562 增加至 9922，状态方程的求解规模扩大了近 4 倍。

然而，如图~\ref{fig:ch4_stop_order_mismatch} 所示，对比 $k=2$ 与 $k=4$ 的优化结果可以清晰地看到：由于控制变量空间的分辨率被刚性锁定，两种阶次下所得的拓扑构型在主承载骨架、杆件特征尺度以及结构宏观边界上高度一致，其差异仅微弱地体现在局部灰度过渡带的插值平滑度上。换言之，单纯提升位移场的逼近阶次，仅仅是对同一幅低分辨率的粗糙几何分布进行冗余的物理场求值，并未从根本上拓宽设计空间的拓扑表达能力。

\begin{figure}[!htbp]
	\centering
	\captionsetup{justification=centering}
	
	\subfigure[$k=2$ ($u_{\mathrm{dof}}=2562$)]{
		\label{fig:ch4_stop_order_k2}
		\includegraphics[width=0.47\textwidth]{ch4_stop_order_k2}
	}
	\hfill
	\subfigure[$k=4$ ($u_{\mathrm{dof}}=9922$)]{
		\label{fig:ch4_stop_order_k4}
		\includegraphics[width=0.47\textwidth]{ch4_stop_order_k4}
	}
	
	\caption{STOP 框架下对称二维 MBB 梁在 $\rho_{\mathrm{dof}}=300$ 时不同有限元阶次的优化拓扑构型对比}
	\label{fig:ch4_stop_order_mismatch}
\end{figure}

进一步地，若考察上述精度错配的对偶情形——即在保持分析网格不变（或有限元阶次较低）的前提下，单方面提高设计空间的分辨率，同样会诱发严重的数值病态。以每个位移分析单元内采用 $4 \times 4$ 的子密度细分为例，此时控制变量的自由度激增至 $\rho_{\text{dof}} = 4800$，而状态方程的求解自由度仍停留在 $u_{\text{dof}} = 2562\ (k=2)$。

\begin{figure}[!htbp]
	\centering
	\captionsetup{justification=centering}
	\includegraphics[width=0.75\textwidth]{ch4_stop_design_refine}
	\caption{分析自由度固定（$u_{\mathrm{dof}}=2562$）时单方面提高设计分辨率（$4\times4$ 子密度）诱发的拓扑数值震荡与伪结构}
	\label{fig:ch4_stop_design_refine}
\end{figure}

如图~\ref{fig:ch4_stop_design_refine} 所示，当分辨率受限的分析模型对设计空间中的高频密度变化存在明显的“欠解析”时，精细的局部拓扑特征无法在有限元分析中得到准确评估。此时，设计空间中不断增加的高频自由度无法被分析模型充分辨识，从而使优化迭代缺乏稳定可靠的物理判据，进而表现为剧烈的密度震荡、条纹状伪结构以及收敛不稳定等现象。

综上所述，精度错配的对偶现象表明：无论单方面提高分析阶次，还是单方面提升设计分辨率，都难以实质性改善优化构型。为了在可控计算代价下充分发挥高阶有限元的逼近优势，并同步提升最优设计的拓扑表达能力，必须打破设计模型与分析模型在离散空间中的刚性耦合。这一底层逻辑直接导向了 “设计--积分--分析” 相互协同的三层多分辨率计算框架，并构成了本章后续建立多分辨率离散策略与密度映射机制的根本动机。


\section{多分辨率拓扑优化计算框架}
\label{sec:mtop_framework}

为实现设计空间与状态空间在离散尺度上的解耦，本章引入并拓展了 Nguyen 等人~\cite{nguyenComputationalParadigmMultiresolution2010a} 提出的 MTOP 范式。该框架打破了传统分片常数密度的离散假定，构建了位移有限元网格、密度积分网格与设计变量网格三套相互独立且逻辑耦合的离散层级。在高阶有限元分析背景下，这种将设计模型与分析模型在离散空间上相对解耦的策略，不仅能够在粗位移网格上实现高分辨率的拓扑表征，也为缓解传统单分辨率框架中的精度错配问题提供了合理的数值途径
%这种将控制变量与状态变量的离散空间予以解耦的策略，不仅能在粗位移网格上实现极高分辨率的拓扑表征，更是彻底规避“精度错配”的必然数学选择
~\cite{nguyenTopologyOptimizationUsing2017a}。


\subsection{多分辨率空间离散策略}
\label{subsec:multires_spatial_discretization}

为了打破传统方法中几何分辨率受限于分析单元尺寸的限制，本章引入了三套相互独立但逻辑关联的离散网格系统，以在数学上实现设计与分析的解耦，并为粗网格高阶分析下的准确数值计算提供支撑~\cite{nguyenComputationalParadigmMultiresolution2010a}。需要说明的是，延续第~\ref{subsec:model_definition_binarity} 节的记号，$\rho(\boldsymbol{x})$ 表示连续意义下的物理密度场，而 $\boldsymbol{\rho}$ 表示该物理密度场在密度积分网格 $\mathcal{T}_\rho$ 上的离散向量表示，其分量由各密度子单元上的物理密度组成。首先是位移有限元网格 $\mathcal{T}_h$，作为求解参数化线弹性平衡方程
\[
\boldsymbol{K}(\boldsymbol{\rho})\boldsymbol{U} = \boldsymbol{F}
\]
的分析模型。在本章研究中，$\mathcal{T}_h$ 由相对较粗的拉格朗日有限元（阶次 $k\geq1$）构成，旨在自由度规模可控的前提下，提高位移场、应变能及其相关约束泛函的近似精度，从而降低由粗网格离散带来的偏微分方程求解误差。

与分析模型相对应的是用于承载优化变量的设计变量网格 $\mathcal{T}_d$。该网格由高分辨率的离散节点组成，每个网格点关联一个设计变量 $d_i \in [0, 1]$。$\mathcal{T}_d$ 的分辨率显著高于 $\mathcal{T}_h$，且二者在几何拓扑上相互独立。这种离散策略允许优化算法在相对粗糙的高阶分析背景下，探索超越分析单元特征尺度的精细拓扑结构，从而显著提高设计空间的拓扑表达能力。

然而，当粗位移单元内部的材料分布受控于极高分辨率的设计变量时，连续物理密度场 $\rho(\boldsymbol{x})$
%物理密度场 $\boldsymbol{\rho}$ 
必然呈现出显著的空间非均匀性。此时，若依然沿用传统的分片常数密度假定，或直接在粗单元上采用常规高斯积分进行刚度评估，将无法准确捕捉材料高频分布对局部算子刚度的真实贡献，进而引发严重的数值积分误差与“欠解析”病态。为此，本章进一步引入了密度积分网格 $\mathcal{T}_\rho$ 作为连接高分辨率设计与低分辨率分析的数值桥梁：对于每个位移分析单元 $\Omega_e$，将其在内部精细划分为若干个密度子单元；通过在子单元层级上对物理材料参数进行离散表达与刚度数值积分，从而确保高度非均匀密度场下有限元算子组装的数学准确性与数值一致性。

值得注意的是，为兼顾高分辨率拓扑表达与偏微分方程求解精度，设计网格 $\mathcal{T}_d$ 与积分网格 $\mathcal{T}_\rho$ 的空间分辨率通常显著高于位移网格 $\mathcal{T}_h$。图~\ref{fig:ch4_mtop_three_meshes} 以双二次四边形单元（$Q_2$）为例，直观阐释了该三层体系的空间拓扑与数值映射流程：在高分辨率下承载设计自由度的变量场 $\boldsymbol{d}$（图 4.3(c)），首先通过正则化映射 $\boldsymbol{\rho} = f(\boldsymbol{d})$ 转换为定义在密度积分网格 $\mathcal{T}_\rho$（图 4.3(b)）上的物理密度场；随后，基于各子单元的数值求积结果，最终组装形成定义在位移网格 $\mathcal{T}_h$（图 4.3(a)）上的全局刚度算子 $\boldsymbol{K}(\boldsymbol{\rho})$。需补充说明的是，图中为简明起见展示了 $\mathcal{T}_d$ 与 $\mathcal{T}_\rho$ 空间分辨率一致（即一一对应）的特例，而在该框架的一般数学定义中，二者的离散尺度可完全独立配置。

\begin{figure}[!htbp]
	\centering
	\captionsetup{justification=centering}
	\includegraphics[width=0.98\textwidth]{ch4_mtop_three_meshes}
	\caption{多分辨率计算框架下的三层网格拓扑及其数值映射关系}
	\label{fig:ch4_mtop_three_meshes}
\end{figure}


\subsection{多分辨率框架下的优化问题列式}
\label{subsec:multires_optimization_formulation}

在上述三层离散网格体系下，拓扑优化问题的数学列式需进行相应的解耦与重构。设 $\boldsymbol{d}$ 为定义在设计变量网格 $\mathcal{T}_d$ 上的设计变量向量，$\boldsymbol{U}$ 为定义在位移有限元网格 $\mathcal{T}_h$ 上的全局节点位移向量。以第~\ref{subsec:min_compliance} 节体积分数约束的柔顺度最小化问题为例，变密度框架下的多分辨率拓扑优化一般离散列式可严格描述如下：
\begin{equation} \label{eq:ch4_mtop_formulation}
	\begin{aligned} 
		\min_{\boldsymbol{d}} \quad & c(\boldsymbol{d}) = \boldsymbol{F}^{\top} \boldsymbol{U} \\ 
		\text{subject to} \quad & \boldsymbol{K}(\boldsymbol{\rho}(\boldsymbol{d})) \boldsymbol{U}(\boldsymbol{d}) = \boldsymbol{F},\\ 
		& g_V(\boldsymbol{d}) = \frac{V(\boldsymbol{\rho}(\boldsymbol{d}))}{V_0} - V_f \leq 0,\\ 
		& d_i \in [0, 1], \quad i = 1, \dots, N_d, 
	\end{aligned}
\end{equation}
其中 $c(\boldsymbol{d})$ 为结构的柔顺度目标函数；$\boldsymbol{K}(\boldsymbol{\rho}(\boldsymbol{d}))$ 为全局刚度矩阵，其构造依赖于由设计变量 $\boldsymbol{d}$ 映射生成的离散物理密度向量 $\boldsymbol{\rho}(\boldsymbol{d})$；$\boldsymbol{F}$ 为全局外载荷向量；$V_0$ 与 $V_f$ 分别为设计域的初始总体积与给定的体积分数上限。

需要说明的是，上式中 $N_d$ 为设计自由度。在多分辨率框架下，结构实际体积 $V(\boldsymbol{\rho}(\boldsymbol{d}))$ 需在积分网格 $\mathcal{T}_\rho$ 上精确计算，即 $V(\boldsymbol{\rho}(\boldsymbol{d})) = \sum_{j=1}^{N_\rho} \rho_j(\boldsymbol{d}) v_j$，其中 $N_\rho$、$v_j$ 与 $\rho_j(\boldsymbol{d})$ 分别为密度子单元的总数、体积与物理密度。此外，尽管纯数学意义下的设计变量 $d_i$ 允许取至下界 $0$，但为避免刚度矩阵奇异，实际物理密度的非零下界（$\rho_{\min} > 0$）将在 $\boldsymbol{d} \to \boldsymbol{\rho}$ 的映射环节予以保证。

与传统单分辨率方法 (第 \ref{chap:lagrange_comparison} 章) 的主要区别在于：多分辨率框架下的离散物理密度向量 $\boldsymbol{\rho}$ 不再与设计变量 $\boldsymbol{d}$ 等价，而是通过广义映射关系 $\boldsymbol{\rho}=f(\boldsymbol{d})$ 建立依赖。该映射算子 $f(\cdot)$ 不仅囊括了设计自由度向积分网格的形函数插值，还统一涵盖了密度过滤与非线性投影等正则化机制，从而在底层逻辑上实现了分析与设计分辨率的独立。需要特别辨析的是，若优化中采用灵敏度过滤，其本质是对梯度信息的正则化处理，而不改变进入刚度组装的物理密度场，因此该映射表现为恒等映射 (即 $\boldsymbol{\rho} = \boldsymbol{d}$)，其过滤效应直接作用于后续的灵敏度更新阶段。在此基础配置下，可直接令子单元密度 $\rho_j = d_j$，若进一步假定设计网格与积分网格完全重合（$\mathcal{T}_d \equiv \mathcal{T}_\rho$）并采用子单元分片常值离散，该恒等映射则可具体表达为局部单元坐标下的形式 $\rho_{e,i} = d_{e,i}$ (其中角标 $e,i$ 分别表征第 $e$ 个宏观位移单元内的第 $i$ 个密度子单元)。


\subsection{基于子单元的刚度矩阵数值积分}
\label{subsec:subcell_stiffness_integration}

在 MTOP 框架下，映射生成的物理密度场 $\rho(\boldsymbol{x})$ 在单个较粗的位移分析单元 $\Omega_e$ 内通常呈现出高度非均匀的空间分布。此时，宏观位移单元的刚度矩阵定义为连续域上的积分形式：
\[
\boldsymbol{K}_e=\int_{\Omega_e}\boldsymbol{B}^{\top}(\boldsymbol{x})\,\boldsymbol{D}(\rho(\boldsymbol{x}))\,\boldsymbol{B}(\boldsymbol{x})\,\mathrm{d}\boldsymbol{x},
\]
其中 $\boldsymbol{B}$ 为应变–位移矩阵，$\boldsymbol{D}(\rho)$ 为由材料插值模型确定的本构矩阵。

由于 $\rho(\boldsymbol{x})$ 及其惩罚项在宏观单元内的剧烈变化，导致上述被积函数呈现出极强的空间非线性，常规的固定阶次高斯求积极易失效并引发严重的数值计算误差。为规避这一高阶分析中的理论瓶颈，本章参考 Nguyen 等~\cite{nguyenTopologyOptimizationUsing2017a} 在 $p$ 型多分辨率框架下的离散策略，将宏观位移单元 $\Omega_e$ 精细划分为 $N_e$ 个独立的密度子单元 $\{\Omega_{e,i}\}_{i=1}^{N_e}$。通过在每个子单元内采用分片常数假定（即密度近似为离散常值 $\rho_{e,i}$），高度非线性的本构矩阵 $\boldsymbol{D}(\rho_{e,i})$ 在子单元局部退化为常数矩阵。这一空间离散降维策略使得整体刚度积分可严密地转换为各子单元上的分段求和：
\[
\boldsymbol{K}_e =\sum_{i=1}^{N_e}\int_{\Omega_{e,i}}\boldsymbol{B}^{\top}(\boldsymbol{x})\,\boldsymbol{D}(\rho_{e,i})\,\boldsymbol{B}(\boldsymbol{x})\,\mathrm d\boldsymbol{x}.
\]

为便于数值实现，上述积分需通过等参映射转移至位移单元的参考域 $\hat{\Omega}_e$ 中进行。记等参映射为 $\boldsymbol{x}=\boldsymbol{x}_e(\boldsymbol\xi)$（其中 $\boldsymbol\xi\in\hat{\Omega}_e$），其雅可比矩阵定义为 $\boldsymbol{J}_e(\boldsymbol\xi)=\frac{\partial\boldsymbol{x}_e(\boldsymbol\xi)}{\partial\boldsymbol\xi}$。进一步地，将参考域 $\hat{\Omega}_e$ 规则划分为 $N_e$ 个参考密度子单元 $\{\hat{\Omega}_{e,i}\}_{i=1}^{N_e}$。为在各个子单元上应用标准高斯求积法则，需引入从标准积分域到第 $i$ 个参考密度子单元的局部映射。设该局部映射的雅可比矩阵为 $\boldsymbol{J}_{\text{sub},i}$，则定义相应的几何缩放因子为
\[
J_i:=\big|\det(\boldsymbol{J}_{\text{sub},i})\big|.
\]
特别地，当参考密度子单元采用规则等分划分时，该缩放因子 $J_i$ 退化为一常数。 对每个参考密度子单元采用 $N_g$ 点高斯积分，设标准积分域上的第 $g$ 个高斯积分点及其对应权重分别为 $\boldsymbol\xi_{g}$ 与 $w_g$，映射至第 $i$ 个子单元后的等参坐标记为 $\boldsymbol\xi_{i,g}$。此时，单元刚度矩阵可由以下离散求和公式高精度近似：
\begin{equation} \label{eq:ch4_subcell_stiffness}
	\boldsymbol{K}_e \approx \sum_{i=1}^{N_e}\sum_{g=1}^{N_g} \boldsymbol{B}^{\top}(\boldsymbol\xi_{i,g})\,\boldsymbol{D}(\rho_{e,i})\,\boldsymbol{B}(\boldsymbol\xi_{i,g})\, \big|\det(\boldsymbol{J}_e(\boldsymbol\xi_{i,g}))\big|\, J_i\, w_g.
\end{equation}
上述分段积分策略的数学本质是以局部空间细化替代全局高阶求积，从而有效规避了强非线性密度场在粗网格上引发的积分震荡与病态。由于密度积分网格与位移插值空间严格解耦，其网格细化无需重构有限元基函数，使得底层实现相对简洁。此外，各子单元的高斯求积在代数操作上彼此独立，天然适合并行的现代计算架构，为大规模高阶拓扑优化奠定了算法基础。


\subsection{多分辨率离散体系下的灵敏度分析}
\label{subsec:sensitivity_analysis}

本节在第 \ref{subsec:multires_optimization_formulation} 节给出的 MTOP 列式及映射关系 $\boldsymbol\rho=f(\boldsymbol{d})$ 的基础上，推导柔顺度目标关于设计变量的解析灵敏度。记 $\rho_{e,i}$ 为位移单元 $e$ 内第 $i$ 个密度子单元的物理密度，$\boldsymbol U_e$ 为全局位移向量 $\boldsymbol U$ 在该单元自由度上的提取向量。为凸显多分辨率框架的链式传导逻辑，以下推导首先给出目标函数对 $\rho_{e,i}$ 的物理灵敏度，再通过映射将其转换为对设计网格 $\mathcal{T}_d$ 上设计变量 $d_{j}$ 的梯度。

对于位移单元 $e$ 内的第 $i$ 个物理密度子单元 $\rho_{e,i}$，由标准伴随变量法可得其柔顺度灵敏度为
\[
\frac{\partial c}{\partial \rho_{e,i}} = -\,\boldsymbol U_e^{\top}\,\frac{\partial \boldsymbol K_e}{\partial \rho_{e,i}}\,\boldsymbol U_e.
\]
由于 $\rho_{e,i}$ 仅在局部影响宏观单元刚度矩阵中对应子单元的积分贡献，结合第~\ref{subsec:subcell_stiffness_integration} 节的子单元离散求积公式，偏导数 $\partial\boldsymbol{K}_e/\partial\rho_{e,i}$ 可表示为该子单元内高斯积分点的累加形式：
\begin{equation} \label{eq:ch4_subcell_stiffness_sens}
	\frac{\partial \boldsymbol{K}_e}{\partial \rho_{e,i}} \approx \sum_{g=1}^{N_g} \boldsymbol{B}^{\top}(\boldsymbol\xi_{i,g})\,\frac{\partial \boldsymbol{D}}{\partial \rho}(\rho_{e,i})\,\boldsymbol{B}(\boldsymbol\xi_{i,g})\, \big|\det(\boldsymbol{J}_e(\boldsymbol\xi_{i,g}))\big|\, J_i\, w_g.
\end{equation}
若采用标准 SIMP 材料插值模型 $\boldsymbol{D}(\rho)=\rho^p\boldsymbol{D}_0$，其中 $\boldsymbol{D}_0$表示实体材料的弹性本构矩阵，则本构矩阵的导数可显式写为
\[
\frac{\partial \boldsymbol{D}}{\partial \rho}(\rho_{e,i}) = p\rho_{e,i}^{p-1}\boldsymbol{D}_0,
\]
将上式代入高斯求积公式中，即可直接通过数值积分高效获取刚度矩阵的局部灵敏度。

进一步，由于物理密度场由映射 $\boldsymbol{\rho}=f(\boldsymbol{d})$ 确定，柔顺度目标函数关于设计变量网格 $\mathcal{T}_d$ 上第 $j$ 个设计变量 $d_j$ 的梯度可基于链式求导法则写为
\[
\frac{\partial c}{\partial d_{j}} = \sum_{e \in \mathcal{N}_j} \sum_{i=1}^{N_e} \frac{\partial c}{\partial \rho_{e,i}} \frac{\partial \rho_{e,i}}{\partial d_{j}},
\]
其中 $\mathcal{N}_j$ 为受设计变量 $d_j$ 影响的位移单元集合。为显式刻画该映射算子的局部紧支撑特性，可定义如下灵敏度贡献集合：
\[
\mathcal{S}_{j} := \left\{(e,i) : \frac{\partial \rho_{e,i}}{\partial d_{j}} \neq 0 \right\},
\]
此时，设计变量的梯度表达式可等价简写为：
\begin{equation} \label{eq:ch4_design_var_grad}
	\frac{\partial c}{\partial d_{j}} = \sum_{(e,i) \in \mathcal{S}_{j}} \frac{\partial c}{\partial \rho_{e,i}} \frac{\partial \rho_{e,i}}{\partial d_{j}}.
\end{equation}
需要指出的是，上述梯度计算依赖于映射算子的雅可比元素 $\partial \rho_{e,i}/\partial d_{j}$。在不引入密度过滤或投影机制且设计变量与子单元满足一一对应关系的简化情形下，该项表现为 Kronecker $\delta$ 形式，即设计变量梯度直接等价于对应子单元的物理灵敏度。而当引入密度过滤或非线性投影等正则化算子后，单个设计变量 $d_{j}$ 的扰动将通过映射权重影响其邻域内的多个物理密度子单元 $\rho_{e,i}$，此时需严格按照集合 $\mathcal{S}_j$ 进行灵敏度贡献的累加。


\section{设计变量到物理密度的映射机制}
\label{sec:mapping}

在第~\ref{sec:mtop_framework} 节建立的三层离散体系中，设计变量 $\boldsymbol{d}$ 定义在设计变量网格 $\mathcal{T}_d$ 上，而刚度矩阵的组装与体积分数计算依赖密度积分网格 $\mathcal{T}_\rho$ 上的离散物理密度向量 $\boldsymbol{\rho}$。为在不同分辨率的网格间建立联系，灵敏度过滤、密度过滤与非线性投影等正则化策略可以统一抽象为设计变量到离散物理密度向量的映射关系 $\boldsymbol{\rho} = f(\boldsymbol{d})$。根据映射关系的数学性质，可将其分为两类：一类为恒等映射，对应灵敏度过滤等梯度正则化方法，即 $\boldsymbol{\rho} = \boldsymbol{d}$，正则化仅作用于梯度场；另一类为密度映射，对应密度过滤与非线性投影等方法，即通过线性或非线性算子对密度场进行重构，使得一个设计变量可影响邻域内多个子单元的物理密度。


\subsection{匹配分辨率下的灵敏度过滤与恒等映射}
\label{subsec:mapping-sens-identity}

当采用灵敏度过滤策略时，正则化操作仅作用于目标函数的梯度场，而不直接改变用于计算的物理密度场。在多分辨率设计中，若设定设计变量网格 $\mathcal{T}_d$ 与密度积分网格 $\mathcal{T}_\rho$ 具有一一对应的拓扑关系（即每个子密度单元由一个独立的设计变量控制），则设计变量到物理密度构成了恒等映射 $\rho_{e,i} = d_j$，其中索引 $j$ 与 $(e,i)$ 在空间位置上一一对应。该映射的分量导数可表示为广义的 Kronecker delta 符号，即 $\partial \rho_{e,i} / \partial d_j = \delta_{(e,i)j}$。根据链式求导法则，目标函数关于设计变量 $d_j$ 的基础梯度即等于对应子密度单元的原始灵敏度：
\[
\frac{\partial c}{\partial d_j} = \sum_{e}\sum_{i}\frac{\partial c}{\partial \rho_{e,i}} \frac{\partial \rho_{e,i}}{\partial d_j} = \frac{\partial c}{\partial \rho_{e,i}}\Bigg|_{(e,i) \leftrightarrow j},
\]
这表明在恒等映射下，第 \ref{subsec:subcell_stiffness_integration} 节基于子单元刚度数值积分得到的 $\partial c/\partial \rho_{e,i}$ 可直接作为设计变量的未过滤梯度。随后，仅需对该原始梯度场施加空间卷积过滤，即可有效抑制棋盘格等数值不稳定性，并向优化算法提供平滑后的下降方向。

参照第~\ref{subsec:cell_density_filtering} 节的离散卷积思想，在密度积分网格 $\mathcal{T}_\rho$ 上定义子密度单元 $(e,i)$ 与其邻域子密度单元 $(t,m)$ 之间的权重：
\[
H_{(e,i)(t,m)} = \max\Bigl\{0,\; r_{\min} - \mathrm{dist}\bigl((e,i), (t,m)\bigr)\Bigr\}^q,
\]
记 $\boldsymbol{x}_{e,i}$ 与 $\boldsymbol{x}_{t,m}$ 分别表示子密度单元 $(e,i)$ 与 $(t,m)$ 的几何中心坐标，并定义两子单元中心之间的距离为
\[
\mathrm{dist}\bigl((e,i), (t,m)\bigr) := \vert\boldsymbol{x}_{e,i} - \boldsymbol{x}_{t,m}\vert,
\]
其中 $q \geq 1$ 为衰减指数（$q=1$ 时退化为线性锥形核）。进一步定义与 $(e,i)$ 相互作用的邻域索引集合：
\[
\mathcal{N}_{e,i} = \{(t,m) \mid H_{(e,i)(t,m)} > 0\},
\]
则灵敏度过滤的离散格式可写为：
\begin{equation} \label{eq:ch4_mtop_sens_filter}
	\widetilde{\left(\frac{\partial c}{\partial \rho_{e,i}}\right)} = \frac{1}{\max\{\gamma, \rho_{e,i}\}}\; \frac{\sum\limits_{(t,m)\in\mathcal{N}_{e,i}} H_{(e,i)(t,m)}\, v_{t,m}\,\rho_{t,m}\, \left(\frac{\partial c}{\partial \rho_{t,m}}\right)} {\sum\limits_{(t,m)\in\mathcal{N}_{e,i}} H_{(e,i)(t,m)}\, v_{t,m}},
\end{equation}
其中，$v_{t,m}$ 为邻域子密度单元 $(t,m)$ 的体积，$\gamma>0$ 为防止除零导致数值奇异的微小正数。由于在恒等映射设定下，设计变量与子单元一一对应（即 $\rho_{e,i} = d_j$），过滤后的 $\widetilde{\partial c/\partial \rho_{e,i}}$ 可直接作为优化更新所需的设计变量 $d_j$ 的梯度使用。该策略实现简单、计算开销低，且无需引入额外的密度映射变量，但由于正则化仅作用于灵敏度而不显式平滑物理密度场。在该设定下，结构边界的几何分辨率主要受子密度单元尺度影响，而最小特征尺度与边界平滑程度则主要由过滤半径 $r_{\min}$ 调节。


\subsection{基于密度过滤的跨网格线性映射}
\label{subsec:mapping_density_linear}

与灵敏度过滤不同，密度过滤直接对密度场进行正则化。在 MTOP 框架下，设计变量 $d_j$ 定义在设计网格 $\mathcal{T}_d$ 上，而物理密度 $\rho_{e,i}$ 定义在密度积分网格 $\mathcal{T}_\rho$（子密度单元）上，因此密度过滤自然对应一类从 $\mathcal{T}_d$ 到 $\mathcal{T}_\rho$ 的跨网格显式映射，该线性映射的离散格式可写为：
\begin{equation} \label{eq:ch4_mtop_density_mapping}
	\rho_{e,i} = \frac{\sum\limits_{j \in \mathcal{N}_{e,i}} H_{(e,i)j}\, v_j\, d_j}{\sum\limits_{j \in \mathcal{N}_{e,i}} H_{(e,i)j}\, v_j}.
\end{equation}
由此可见，密度过滤本质上将每个 $\rho_{e,i}$ 表示为邻域内设计变量 $d_j$ 的加权平均，因此映射算子 $f(\cdot)$ 表现为典型的线性算子。其中，$H_{(e,i)j}$ 为积分点（子单元几何中心）$\boldsymbol{x}_{e,i}$ 与设计变量坐标 $\boldsymbol{x}_j$ 之间的权重函数，$\mathcal{N}_{e,i}$ 为对子单元 $(e,i)$ 有贡献的设计变量索引集合，$v_j$ 为设计变量对应的参考体积。由于 $\rho_{e,i}$ 由邻域内多个 $d_j$ 共同决定，该映射对设计变量的偏导数可精确求得：
\begin{equation} \label{eq:ch4_mtop_density_filter_deriv}
	\frac{\partial \rho_{e,i}}{\partial d_j} = 
	\begin{cases} 
		\dfrac{H_{(e,i)j}\, v_j}{\sum\limits_{m \in \mathcal{N}_{e,i}} H_{(e,i)m}\, v_m}, & j \in \mathcal{N}_{e,i}, \\[1.5em] 
		0, & j \notin \mathcal{N}_{e,i}. 
	\end{cases}
\end{equation}
结合上述映射关系，目标函数 $c$ 关于设计变量 $d_j$ 的总灵敏度可通过链式求导法则完整传导：
\begin{equation} \label{eq:ch4_mtop_density_filter_chain}
	\frac{\partial c}{\partial d_j} = \sum_{(e,i) \in \mathcal{S}_j} \frac{\partial c}{\partial \rho_{e,i}} \frac{\partial \rho_{e,i}}{\partial d_j},
\end{equation}
其中 $\mathcal{S}_j$ 为受设计变量 $d_j$ 影响的所有物理子单元的索引集合。同理，体积分数等结构约束关于设计变量的梯度亦可遵循该链式法则框架进行统一计算。


\subsection{面向高分辨率边界的子单元非线性投影}
\label{subsec:mapping-heaviside-nonlinear}

如第~\ref{subsec:projection_methods} 节所述，虽然密度过滤策略能够有效控制结构的最小特征尺寸并解决棋盘格问题，但其本质上的空间加权平均操作不可避免地会在结构边界处引入较宽的灰度过渡带。在 MTOP 中，为了充分利用高分辨率密度积分网格 $\mathcal{T}_{\rho}$ 来清晰地表达结构边界，本节在密度过滤的基础上进一步引入 Heaviside 近似阶跃函数作为非线性投影算子，建立从“设计变量 $\to$ 中间密度 $\to$ 物理密度”的非线性映射机制。

引入投影算子后，首先利用第~\ref{subsec:mapping_density_linear} 节所述的线性密度过滤，将定义在独立设计网格 $\mathcal{T}_d$ 上的设计变量 $d_{j}$ 映射为定义在密度积分网格 $\mathcal{T}_{\rho}$（子密度单元）上的中间密度 $\tilde{\rho}_{e,i}$。随后，在密度积分网格上对中间密度 $\tilde{\rho}_{e,i}$ 施加非线性投影，从而得到高度逼近 $0$-$1$ 分布的物理密度 $\rho_{e,i}$。常用的光滑投影函数主要包括以下两类。
\begin{description}
	\item[指数型投影]：其离散形式为
	\[
	\rho_{e,i} = 1 - e^{-\beta\tilde{\rho}_{e,i}} + \tilde{\rho}_{e,i}e^{-\beta}.
	\]
	
	\item[双曲正切型投影]：其离散形式为
	\[
	\rho_{e,i} = \frac{\tanh(\beta\eta) + \tanh(\beta(\tilde{\rho}_{e,i} - \eta))}{\tanh(\beta\eta) + \tanh(\beta(1 - \eta))}.
	\]
\end{description}
其中，$\beta$ 为控制投影陡峭程度的正则化参数，$\eta\in(0,1)$ 为投影阈值（实际计算中常取 $\eta=0.5$ 以近似保持投影前后的体积守恒）。为避免强非线性投影导致优化陷入不利的局部极值，通常对 $\beta$ 采用延拓（Continuation）策略，在迭代中逐步增大其值，从而稳健地驱动灰度单元向清晰的物理边界演化。值得注意的是，在 MTOP 框架下，该非线性投影直接在密度积分网格 $\mathcal{T}_{\rho}$ 上逐点执行，这意味着即使设计网格 $\mathcal{T}_d$ 上的变量分布相对稀疏，依然能通过在极细密的积分网格上施加锐利的非线性截断，获得边界更清晰、二值性更强的高分辨率拓扑表达。

引入投影机制后，目标函数 $c$ 关于设计变量 $d_j$ 的灵敏度计算需严格应用链式求导法则。结合第 \ref{subsec:mapping_density_linear} 节中线性过滤的推导，总灵敏度公式扩展为
\[
\frac{\partial c}{\partial d_j} = \sum_{(e,i) \in \mathcal{S}_j} \frac{\partial c}{\partial \rho_{e,i}} \frac{\partial \rho_{e,i}}{\partial \tilde{\rho}_{e,i}} \frac{\partial \tilde{\rho}_{e,i}}{\partial d_j},
\]
其中，$\frac{\partial c}{\partial \rho_{e,i}}$ 为目标函数对子单元物理密度的偏导数；$\frac{\partial \tilde{\rho}_{e,i}}{\partial d_j}$ 为线性密度过滤的权重项；$\frac{\partial \rho_{e,i}}{\partial \tilde{\rho}_{e,i}}$ 为非线性投影函数的导数项。

\begin{description}
	\item[指数型投影的导数项]：
	\[
	\frac{\partial \rho_{e,i}}{\partial \tilde{\rho}_{e,i}} = \beta e^{-\beta\tilde{\rho}_{e,i}} + e^{-\beta}.
	\]
	
	\item[双曲正切型投影的导数项]：
	\[
	\frac{\partial \rho_{e,i}}{\partial \tilde{\rho}_{e,i}} = \beta \frac{1 - \tanh^2(\beta(\tilde{\rho}_{e,i} - \eta))}{\tanh(\beta\eta) + \tanh(\beta(1 - \eta))}.
	\]
\end{description}
通过上述完备的链式法则，优化驱动力（灵敏度信息）能够准确无误地穿透非线性投影层与线性过滤层，最终传递回独立的设计变量网格上，从而稳健地驱动拓扑构型的演化。


\section{多分辨率应力约束拓扑优化模型与算法}
\label{sec:ch4_mtop_stress_constraint}

在连续体结构拓扑优化中，局部应力约束设计长期面临着应力奇异性与大规模局部约束两大数值挑战（如第~\ref{subsec:stress_constrained} 节所述）。为此，本节引入多项式消失约束技术以有效消除低密度区域的非物理应力奇异性，并采用 ALM 来稳健处理大规模且相互独立的局部约束。在此基础上，本节将上述策略与多分辨率离散架构深度融合，构建面向高阶有限元的多分辨率应力约束拓扑优化模型，并详细推导其相应的数值求解算法与伴随灵敏度过程。


\subsection{应力约束优化模型列式}
\label{subsec:stress_constraint_model}

考虑结构轻量化设计中最具代表性的问题，即应力约束下的体积最小化问题。该问题的核心在于，在满足线弹性平衡方程且全域材料点局部应力均不超过许用值的前提下，寻求最优的材料空间分布，从而使结构的总体积最小。在 MTOP 框架下，设计变量与应力约束被分别定义在具有不同分辨率的离散网格上。具体而言，设计变量 $d_j$ 定义在独立的设计网格 $\mathcal{T}_d$ 上，而局部应力的评估则下沉至位移单元内部的高分辨率密度积分网格 $\mathcal{T}_\rho$（即子密度单元）层面。该问题的原始数学模型可严格表述为：
\begin{equation} \label{eq:ch4_stress_mtop_model}
	\begin{aligned}
		\min_{\boldsymbol{d}} \quad & f(\boldsymbol{d}) = \frac{V(\boldsymbol{\rho}(\boldsymbol{d}))}{V_0} \\
		\text{subject to} \quad & \boldsymbol{K}(\boldsymbol{\rho}(\boldsymbol{d}))\boldsymbol{U}(\boldsymbol{d}) = \boldsymbol{F}, \\
		& \hat{g}_{e,i}(\boldsymbol{d}, \boldsymbol{U}) = \sigma_{e,i}^{\text{vM}}(\boldsymbol{U}, \rho_{e,i}) - \bar{\sigma} \le 0, \quad \forall e \in \mathcal{T}_h, \, i = 1, \dots, N_e, \\
		& 0 \le d_j \le 1, \quad j = 1, \dots, N_d.
	\end{aligned}
\end{equation}
其中，$f(\boldsymbol{d})$ 为目标函数（体积分数）；$N_d$ 表示设计变量的总数，$N_e$ 表示每个位移分析单元内包含的密度子单元数量；$d_j$ 表示设计网格 $\mathcal{T}_d$ 上的第 $j$ 个设计变量；$\sigma_{e,i}^{\text{vM}}$ 表示第 $e$ 个位移单元内第 $i$ 个密度子单元几何中心处的 von Mises 等效应力，$\bar{\sigma}$ 为材料的许用应力。

为了充分发挥高阶有限元的高精度特性，局部应力场的评估直接基于高阶形函数梯度，在精细的子单元几何中心处进行精准采样。相较于传统低阶单元的分片常应力假定，该评估机制能够敏锐地捕捉到结构内部复杂的应力集中与局部高梯度响应，从而提高局部应力约束判定的精度与应力分析结果的可靠性。


\subsection{增广拉格朗日应力约束框架}
\label{subsec:alm_stress_constraint}

为了在多分辨率架构下实现对高阶局部应力场的精准控制，本节构建了基于 ALM 的求解框架。该框架保留了全域所有的局部约束，并采用多项式消失约束技术来彻底克服低密度区域的数值病态。在 MTOP 框架中，应力约束被直接施加在极精细的密度积分网格 $\mathcal{T}_\rho$ 上。为解决优化过程中的应力奇异性难题，本节采用 Giraldo-Londoño 和 Paulino~\cite{giraldo-londonoUnifiedApproachTopology2020} 提出的多项式消失约束技术对原始屈服准则进行数学松弛。首先，定义第 $e$ 个位移分析单元内第 $i$ 个密度子单元的归一化应力因子：
\[
\Lambda_{e,i} = \frac{\sigma_{e,i}^{\text{vM}}}{\bar{\sigma}} - 1,
\]
基于此，将原始的不等式约束  $\hat{g}_{e,i}(\boldsymbol{d}, \boldsymbol{U}) \le 0$ 松弛为如下形式的连续函数：
\begin{equation} 
	g_{e,i}(\boldsymbol{d}, \boldsymbol{U}) = m_E(\rho_{e,i})\Lambda_{e,i}(\Lambda_{e,i}^2 + 1),
	\label{eq:ch4_poly_restrict_func}
\end{equation}
其中，$m_E(\rho_{e,i})$ 为与刚度惩罚模型保持一致的材料插值函数。该松弛策略的数学机理十分明确：当 $\rho_{e,i} \to 0$（孔洞区域）时，$m_E(\rho_{e,i}) \to 0$，约束函数恒趋于 0，从而自动屏蔽虚假应力并消除奇异性；而当 $\rho_{e,i} \approx 1$（实体区域）时，$m_E(\rho_{e,i}) \approx 1$，松弛约束退化为 $\Lambda_{e,i}(\Lambda_{e,i}^2 + 1) \le 0$，由于 $\Lambda_{e,i}^2 + 1 > 0$，上述松弛约束严格等价于 $\Lambda_{e,i} \le 0$，恢复了真实的物理屈服准则 $\sigma_{e,i}^{\text{vM}} \le \bar{\sigma}$。

在 ALM 中，原始的带约束优化问题被转化为求解一系列无约束子问题。特别地，在多分辨率架构下，大规模局部约束往往会导致惩罚项的量级过大，进而掩盖真实目标函数（体积分数）的优化驱动力。作为缓解此问题的一种有效手段，Senhora 等~\cite{senhoraTopologyOptimizationLocal2020} 在增广拉格朗日惩罚项 $P^{(k)}(\boldsymbol{d}, \boldsymbol{U})$ 中引入了一个归一化因子。在第 $k$ 个增广拉格朗日迭代步中，优化的归一化无约束子问题可写为：
\begin{equation} \label{eq:ch4_alm_subproblem}
	\min_{\boldsymbol{d}}~\mathcal{J}^{(k)}(\boldsymbol{d}, \boldsymbol{U}) = f(\boldsymbol{d}) + \frac{1}{N_\rho} P^{(k)}(\boldsymbol{d}, \boldsymbol{U}),
\end{equation}
其中，$N_\rho$ 为密度积分网格上的局部应力约束总数，$P^{(k)}(\boldsymbol{d}, \boldsymbol{U})$ 为增广拉格朗日惩罚项。为了严格映射多分辨率网格的层级结构，该惩罚项展开为对位移单元及其内部子单元的双重求和形式：
\[
P^{(k)}(\boldsymbol{d}, \boldsymbol{U}) = \sum_{e \in \mathcal{T}_h} \sum_{i=1}^{N_e} \left[ \lambda_{e,i}^{(k)} h_{e,i}(\boldsymbol{d}, \boldsymbol{U}) + \frac{\mu^{(k)}}{2} h_{e,i}(\boldsymbol{d}, \boldsymbol{U})^2 \right],
\]
式中，$\lambda_{e,i}^{(k)}$ 为第 $k$ 步迭代中对应于第 $e$ 个位移单元内第 $i$ 个密度子单元局部约束的拉格朗日乘子估计值；$\mu^{(k)} > 0$ 为二次惩罚因子；$h_{e,i}(\boldsymbol{d}, \boldsymbol{U})$ 为结合了不等式约束特性的等效约束函数，其表达式为
\[
h_{e,i}(\boldsymbol{d}, \boldsymbol{U}) = \max \left[ g_{e,i}(\boldsymbol{d}, \boldsymbol{U}), -\frac{\lambda_{e,i}^{(k)}}{\mu^{(k)}} \right].
\]
通过引入系数 $1/N_\rho$ 对惩罚项进行归一化处理，可以有效保证在约束数量随高分辨率网格剧增时，目标函数与惩罚项在数值演化过程中始终保持合理的量级平衡，从而显著提升算法在大规模局部应力约束问题上的稳健性与收敛效率。

为了直观地展示上述求解过程，图~\ref{fig:alm_flowchart} 给出了本研究所实现的基于归一化 ALM 框架的拓扑优化算法流程图。
\begin{figure}[H]
	\centering
	\begin{tikzpicture}[
		scale=0.85, transform shape,
		node distance=0.8cm and 2cm,
		% 1. 基础扁平化样式
		base/.style = {draw=black, align=center, minimum height=1.2cm},
		% 为 io 添加了 inner xsep 以控制文字与边缘的距离，通过内部换行来控制总宽度
		io/.style = {base, trapezium, trapezium left angle=70, trapezium right angle=110, fill=gray!10, thick, inner xsep=0.6cm},
		process/.style = {base, rectangle, fill=white, minimum width=6cm, thick},
		% 2. 核心优化模块
		mainopt/.style = {base, rectangle, fill=blue!5, minimum width=6.5cm, minimum height=1.6cm, draw=blue!40!black, very thick}, 
		decision/.style = {base, diamond, aspect=2.5, fill=yellow!10, minimum width=3cm, thick},
		% 3. 虚线模块框
		groupbox/.style = {draw=black!50, dashed, thick, rounded corners=0.2cm, inner sep=0.4cm},
		% 4. 箭头样式
		arrow/.style={thick, -{Stealth[scale=1.2]}}
		]
		
		% --- 放置主要节点 ---
		\node (in) [io] {\heiti 读取初始数据};
		\node (init) [process, below=of in] {{\heiti 拓扑优化参数初始化}\\ $k=0, \lambda_{e,i}^{(0)}=0, \mu^{(0)}>0$};
		
		\node (mma) [mainopt, below=0.9cm of init] {{\heiti 求解归一化 ALM 无约束子问题 (MMA)}\\ $\displaystyle\mathcal{J}^{(k)}(\boldsymbol{d}, \boldsymbol{U}) = f(\boldsymbol{d}) + \frac{1}{N_\rho} P^{(k)}(\boldsymbol{d}, \boldsymbol{U})$};
		
		\node (lam) [process, below=1.2cm of mma] {{\heiti 拉格朗日乘子更新}\\ $\displaystyle \lambda_{e,i}^{(k+1)} = \lambda_{e,i}^{(k)} + \mu^{(k)}h_{e,i}(\boldsymbol{d}^{(k)}, \boldsymbol{U}^{(k)})$};
		\node (mu) [process, below=0.6cm of lam] {{\heiti 二次惩罚因子更新}\\ $\displaystyle \mu^{(k+1)} = \min\big[\alpha\mu^{(k)}, \mu_{\max}\big]$};
		
		% --- 绘制虚线包裹框 ---
		\node (updategroup) [groupbox, fit=(lam) (mu)] {};
		\node [fill=white, anchor=center, inner sep=3pt, text=black!70] at (updategroup.north) {\sffamily \small ALM 乘子与惩罚项更新模块};
		
		% 【修改 2】：修改为满足应力收敛准则
		\node (conv) [decision, below=0.9cm of updategroup] {\heiti 满足应力收敛准则？};
		% 【修改 3】：精简输出文本
		\node (out) [io, below=of conv] {\heiti 输出最优拓扑构型};
		
		% --- 绘制主流程连接箭头 ---
		\draw [arrow] (in) -- (init);
		\draw [arrow] (init) -- (mma);
		\draw [arrow] (mma) -- (updategroup.north); 
		\draw [arrow] (lam) -- (mu);
		\draw [arrow] (updategroup.south) -- (conv); 
		\draw [arrow] (conv) -- node[anchor=west] {是} (out);
		
		% --- 绘制循环返回箭头 ---
		\draw [arrow] (conv.east) -- node[anchor=south] {否} ++(3,0) |- node[anchor=west, pos=0.25] {$k \leftarrow k+1$} (mma.east);
		
	\end{tikzpicture}
	
	\vspace{0.2cm}
	\caption{基于 ALM 局部应力约束框架的拓扑优化算法流程图}
	\label{fig:alm_flowchart}
\end{figure}

在完成归一化子问题 $\mathcal{J}^{(k)}(\boldsymbol{d}, \boldsymbol{U})$ 的近似求解后，需对拉格朗日乘子 $\lambda_{e,i}$ 与二次惩罚因子 $\mu$ 进行更新，以驱动算法向可行域逼近：
\begin{equation} \label{eq:ch4_multiplier_update}
	\lambda_{e,i}^{(k+1)} = \lambda_{e,i}^{(k)} + \mu^{(k)} h_{e,i}(\boldsymbol{d}^{(k)}, \boldsymbol{U}^{(k)}),
\end{equation}
\begin{equation} \label{eq:ch4_penalty_update}
	\mu^{(k+1)} = \min \left[ \alpha \mu^{(k)}, \mu_{\max} \right].
\end{equation}
其中，$\alpha > 1$ 为惩罚因子的更新参数，$\mu_{\max}$ 为防止数值病态而设置的惩罚因子上限。在优化算法初始阶段，通常取 $\lambda_{e,i}^{(0)} = 0$，并将 $\mu^{(0)}$ 设为一较小的正数。

ALM 方法的收敛性能对超参数的设置较为敏感。结合数值实现经验，各参数的选取原则归纳如下~\cite{giraldo-londonoPolyStressMatlabImplementation2021}：
\begin{enumerate}
	\item 初始拉格朗日乘子 $\lambda_{e,i}^{(0)}$：对优化初期的性能影响最小，通常可直接初始化为零向量。
	
	\item 初始惩罚因子 $\mu^{(0)}$：需兼顾约束满足的演化速度。若其值过小，会导致算法向可行域收敛的过程过于缓慢，通常需通过初始校准获取一个能与目标函数量级匹配的量值。
	
	\item 惩罚因子上限 $\mu_{\max}$：主要用于保障数值稳定性。若该值过大，会导致全局刚度矩阵的条件数在优化后期急剧恶化，进而引发数值病态。
	
	\item 更新参数 $\alpha$：建议选取较小的值（例如 $1 < \alpha < 2$），以确保惩罚因子 $\mu^{(k)}$ 在连续子问题之间温和增长，从而保障优化轨迹平滑性的稳定性。
\end{enumerate}
数值实验表明，一旦确定了适用于当前离散尺度的超参数组合，该配置通常能够稳健地推广至各类复杂几何拓扑问题的求解中。


\subsection{增广拉格朗日函数的灵敏度分析}
\label{subsec:al_sensitivity}

为采用基于梯度的优化算法求解上述应力约束拓扑优化问题，必须提供归一化增广拉格朗日无约束子问题对设计变量的准确梯度信息。在本框架中，由于目标函数与惩罚项相互独立，根据链式法则，目标泛函 $\mathcal{J}^{(k)}(\boldsymbol{d}, \boldsymbol{U})$ 关于设计网格 $\mathcal{T}_d$ 上任意设计变量 $d_j$ 的总灵敏度可以写为
\[
\frac{\mathrm{d}\mathcal{J}^{(k)}}{\mathrm{d}d_j} = \frac{\partial f}{\partial d_j} + \frac{1}{N_\rho} \frac{\mathrm{d}P^{(k)}}{\mathrm{d}d_j},
\]
其中，体积分数目标函数 $f(\boldsymbol{d})$ 对设计变量的导数 $\partial f/\partial d_j$ 具有显式的表达形式且计算简便。本节的重点在于局部应力惩罚项 $P^{(k)}$ 对设计变量全导数 $\mathrm{d}P^{(k)}/\mathrm{d}d_j$ 的求解。

在多分辨率框架下，由于引入了密度过滤与投影映射，设计变量 $d_j$ 被映射为密度积分网格 $\mathcal{T}_\rho$ 上的子单元物理密度 $\rho_{e,i}$。因此，惩罚项 $P^{(k)}$ 关于设计变量的导数可通过物理密度场作为中间变量进行链式展开：
\[
\frac{\mathrm{d}P^{(k)}}{\mathrm{d}d_j} = \sum_{(e,i) \in \mathcal{S}_j} \frac{\partial \rho_{e,i}}{\partial d_j} \frac{\mathrm{d}P^{(k)}}{\mathrm{d}\rho_{e,i}}
\]
式中，$\mathcal{S}_j$ 表示受第 $j$ 个设计变量影响的密度子单元索引集合；$\partial \rho_{e,i}/\partial d_j$ 仅由预定义的过滤函数与投影算子决定；$\mathrm{d}P^{(k)}/\mathrm{d}\rho_{e,i}$ 的求解需借助于伴随变量法来消除隐式状态方程的约束。

惩罚项 $P^{(k)}$ 既通过材料插值函数显式依赖于物理密度 $\rho_{e,i}$，又通过位移场 $U$ 隐式依赖于 $\rho_{e,i}$。因此，其对 $\rho_{e,i}$ 的全导数为
\[
\frac{\mathrm{d}P^{(k)}}{\mathrm{d}\rho_{e,i}} = \frac{\partial P^{(k)}}{\partial\rho_{e,i}} + \left(\frac{\partial \boldsymbol{U}}{\partial\rho_{e,i}}\right)^{\top}\frac{\partial P^{(k)}}{\partial \boldsymbol{U}}.
\]
直接计算上式中位移向量对局部密度的导数 $\partial \boldsymbol{U}/\partial\rho_{e,i}$ 需要针对每一个密度子单元求解一次大型线性方程组，其计算成本在大规模多分辨率网格下是不可接受的。为此，本文采用伴随变量法来消除这一隐式求导过程。对系统平衡方程 $\boldsymbol{K}(\rho)\boldsymbol{U}=\boldsymbol{F}$ 两边同时对 $\rho_{e,i}$ 求导，可得
\begin{equation}
	\boldsymbol{K}\frac{\partial \boldsymbol{U}}{\partial\rho_{e,i}} + \frac{\partial \boldsymbol{K}}{\partial\rho_{e,i}}\boldsymbol{U} = \boldsymbol{0},
	\label{eq:ch4_state_eq_derivative}
\end{equation}
引入一个任意的伴随向量 $\boldsymbol{\psi}$，并利用 \eqref{eq:ch4_state_eq_derivative} 式，将全导数表达式改写为
\[
\frac{\mathrm{d} P^{(k)}}{\mathrm{d} \rho_{e,i}} = \frac{\partial P^{(k)}}{\partial \rho_{e,i}} + \left( \frac{\partial \boldsymbol{U}}{\partial \rho_{e,i}} \right)^{\!\top} \frac{\partial P^{(k)}}{\partial \boldsymbol{U}} + \boldsymbol{\psi}^{\top} \left( \boldsymbol{K} \frac{\partial \boldsymbol{U}}{\partial \rho_{e,i}} + \frac{\partial \boldsymbol{K}}{\partial \rho_{e,i}} \boldsymbol{U} \right).
\]
考虑到全局刚度矩阵的对称性（$\boldsymbol{K} = \boldsymbol{K}^T$），为了消除包含隐式导数 $\partial \boldsymbol{U}/\partial\rho_{e,i}$ 的项，令其系数为零，从而构造如下伴随方程：
\begin{equation} \label{eq:ch4_adjoint_equation}
	\boldsymbol{K}\boldsymbol{\psi} = -\frac{\partial P^{(k)}}{\partial \boldsymbol{U}}.
\end{equation}
通过求解该伴随方程获得伴随向量 $\boldsymbol{\psi}$ 后，复杂的全导数计算便可简化为仅依赖于显式项的运算：
\[
\frac{\mathrm{d}P^{(k)}}{\mathrm{d}\rho_{e,i}} = \frac{\partial P^{(k)}}{\partial\rho_{e,i}} + \boldsymbol{\psi}^{\top}\frac{\partial \boldsymbol{K}}{\partial\rho_{e,i}}\boldsymbol{U}.
\]

特别需要指出的是，在 MTOP 框架中，上式的伴随灵敏度项 $\boldsymbol{\psi}^T \frac{\partial \boldsymbol{K}}{\partial\rho_{e,i}} \boldsymbol{U}$ 与传统单分辨率算法存在显著的物理映射差异。由于密度场被细化至位移单元内部的子单元层级，全局刚度矩阵对特定子单元密度 $\rho_{e,i}$ 的偏导数实际上仅作用于局部单元刚度矩阵 $\boldsymbol{K}_e$。为避免传统近似带来的尺度失真，本文采用第 \ref{subsec:subcell_stiffness_integration} 节中建立的子单元精确积分体系。局部刚度矩阵对子单元密度的偏导数 $\partial \boldsymbol{K}_e / \partial\rho_{e,i}$ 可通过对该子单元内部的高斯积分点进行累加显式求得：
\[
\frac{\partial \boldsymbol{K}_e}{\partial \rho_{e,i}} \approx \sum_{g=1}^{N_g} \left( \boldsymbol{B}^T \frac{\partial \boldsymbol{D}}{\partial \rho}(\rho_{e,i}) \boldsymbol{B} \right)(\boldsymbol{x}_e(\boldsymbol\xi_{i,g})) |\det(\boldsymbol{J}_e(\boldsymbol\xi_{i,g}))|J_iw_g,
\]
式中，$J_i$ 为从参考位移单元到子单元的几何缩放因子，它严格量化了多分辨率网格间的尺度差异。将该单元级显式导数代入伴随方程求出的全局关系中，多分辨率架构下的罚函数惩罚项关于局部密度的总灵敏度最终严格表述为：
\begin{equation} 
	\frac{\mathrm{d}P^{(k)}}{\mathrm{d}\rho_{e,i}} = \underbrace{\frac{\partial P^{(k)}}{\partial\rho_{e,i}}}_{\text{显式敏度项}} + \underbrace{\boldsymbol{\psi}_e^T \frac{\partial \boldsymbol{K}_e}{\partial\rho_{e,i}} \boldsymbol{U}_e}_{\text{伴随敏度项}},
	\label{eq:ch4_total_sensitivity}
\end{equation}
式中，$\boldsymbol{U}_e$ 与 $\boldsymbol{\psi}_e$ 分别为全局位移向量与伴随向量在第 $e$ 个单元自由度上的提取分量。在求解包含大规模局部约束的高精度问题时，依托上述严谨的子单元积分求导体系，可有效避免隐式梯度的量级错乱，从而保障应力约束优化求解的数值稳定性与收敛性。

为了完整求解伴随方程与灵敏度，还需要明确偏导数 $\partial P^{(k)}/\partial \boldsymbol{U}$ 和 $\partial P^{(k)}/\partial\rho_{e,i}$ 的具体表达式。根据 $P^{(k)}$ 的双重求和定义，其对全局位移向量的偏导数为
\begin{equation}
	\frac{\partial P^{(k)}}{\partial \boldsymbol{U}} = \sum_{e \in \mathcal{T}_h} \sum_{i=1}^{N_e} \left[ \lambda_{e,i}^{(k)} + \mu^{(k)} h_{e,i}(\boldsymbol{d}, \boldsymbol{U}) \right] \frac{\partial h_{e,i}}{\partial \boldsymbol{U}},
	\label{eq:ch4_penalty_deriv_U}
\end{equation}
由于局部等效约束函数 $h_{e,i}$ 在空间上仅显式依赖于其所在的特定子单元物理密度 $\rho_{e,i}$，\eqref{eq:ch4_penalty_deriv_U} 式中的双重求和号在对某一特定子单元密度求偏导时自动消去，从而退化为极其简洁的标量表达式：
\[
\frac{\partial P^{(k)}}{\partial\rho_{e,i}} = \left[ \lambda_{e,i}^{(k)} + \mu^{(k)} h_{e,i}(\boldsymbol{d}, \boldsymbol{U}) \right] \frac{\partial h_{e,i}}{\partial\rho_{e,i}}.
\]
基于等效约束函数 $h_{e,i}(\boldsymbol{d}, \boldsymbol{U})$ 的最大值函数形式，其导数运算呈现典型的分段特征。具体而言，针对任意变量 $(\cdot)$（代表 $\boldsymbol{U}$ 或 $\rho_{e,i}$），其偏导数可严格写为：
\begin{equation} \label{eq:ch4_h_derivative}
	\frac{\partial h_{e,i}}{\partial(\cdot)} = 
	\begin{cases}
		\displaystyle \frac{\partial g_{e,i}}{\partial(\cdot)}, &\quad \text{若 } g_{e,i}(\boldsymbol{d}, \boldsymbol{U}) \ge -\frac{\lambda_{e,i}^{(k)}}{\mu^{(k)}}, \\
		\boldsymbol{0}, &\quad \text{其他情况}.
	\end{cases}
\end{equation}
上式表明：当局部约束满足 $g_{e,i} < -\lambda_{e,i}^{(k)}/\mu^{(k)}$ 时，该约束在当前增广拉格朗日迭代步中处于非激活状态，不再提供任何梯度驱动力；反之，处于激活状态时，其梯度则直接等价于原始约束函数 $g_{e,i}$ 对相应变量的偏导数。

针对处于激活状态的局部应力约束，基于经过多项式松弛的约束函数~\eqref{eq:ch4_poly_restrict_func}，其关于子单元物理密度 $\rho_{e,i}$ 的显式偏导数可直接计算如下：
\[
\frac{\partial g_{e,i}}{\partial\rho_{e,i}} = \frac{\partial m_E(\rho_{e,i})}{\partial\rho_{e,i}}\Lambda_{e,i}(\Lambda_{e,i}^2 + 1).
\]
另一方面，该约束函数关于全局位移场 $\boldsymbol{U}$ 的偏导数，可通过应力张量分量进行链式展开求得：
\[
\frac{\partial g_{e,i}}{\partial \boldsymbol{U}} = \frac{\partial g_{e,i}}{\partial\sigma_{e,i}^{\text{vM}}} \frac{\partial\sigma_{e,i}^{\text{vM}}}{\partial\boldsymbol{\sigma}_{e,i}} \frac{\partial\boldsymbol{\sigma}_{e,i}}{\partial \boldsymbol{U}},
\]
式中，右侧第一项为松弛约束函数关于等效应力标量的偏导数：
\[
\frac{\partial g_{e,i}}{\partial\sigma_{e,i}^{\text{vM}}} = m_E(\rho_{e,i})\frac{3\Lambda_{e,i}^2 + 1}{\bar{\sigma}},
\]
第二项利用 von Mises 应力二次型表达式 $\sigma_{e,i}^{\text{vM}} = \sqrt{\boldsymbol{\sigma}_{e,i}^T \boldsymbol{M} \boldsymbol{\sigma}_{e,i}}$，对其应力向量 $\boldsymbol{\sigma}_{e,i}$ 求偏导可得：
\[
\frac{\partial\sigma_{e,i}^{\text{vM}}}{\partial\boldsymbol{\sigma}_{e,i}} = \frac{\boldsymbol{\sigma}_{e,i}^T \boldsymbol{M}}{\sigma_{e,i}^{\text{vM}}},
\]
第三项表示局部应力向量关于全局位移的敏感度，可通过材料本构关系与几何方程显式给出：
\[
\frac{\partial\boldsymbol{\sigma}_{e,i}}{\partial \boldsymbol{U}} = \boldsymbol{D}_0\boldsymbol{B}_{e,i},
\]
式中，$\boldsymbol{D}_0$ 为实体材料的弹性本构矩阵，$\boldsymbol{B}_{e,i}$ 为对应局部约束点（即第 $e$ 个单元内第 $i$ 个子单元几何中心）处的高阶应变--位移插值矩阵。将上述各项代回原伴随方程与灵敏度公式中，即可完整构建用于多分辨率应力约束拓扑优化的灵敏度解析列式。


\section{数值算例}
\label{sec:ch4_numerical_examples}

本节旨在通过一系列典型数值算例，对前文提出的基于高阶有限元的 MTOP 进行系统验证与综合评估。在统一的计算基准下，本节将依次探讨二维 MBB 梁问题、不同密度映射机制的对比，以及极具挑战性的局部应力约束应用推广。通过严格的配对实验与分析，系统剖析分析阶次、设计分辨率、以及正则化映射策略（恒等映射、线性过滤与非线性投影）对目标函数收敛行为、应力场逼近精度及最终拓扑构型的具体影响机制。


\subsection{算例设置与统一框架}
\label{subsec:ch4_setup_framework}

为保证本章数值算例在统一且可复现的基准下进行，除特别说明外，均沿用第~\ref{subsec:ch3_case_framework} 节中关于小变形线弹性假设、各向同性材料、SIMP 材料插值、单位体系等通用设定。本章在此基础上，针对多分辨率框架的数值特性，主要补充如下统一约定：
\begin{description}
	\item [惩罚因子延续]：与前文固定惩罚因子（$p=3$）不同，本章为增强多分辨率框架下优化过程的鲁棒性并避免局部极值，对惩罚指数 $p$ 施加如下延续策略：
	\[
	p^{(n)}=\min\Bigl(3,\; 1+0.5\,\bigl\lfloor n/30 \bigr\rfloor \Bigr),
	\]
	即初始 $p=1$，每迭代 30 次增加 0.5，直至达到上限 $p=3$ 后保持不变。
	
	\item [映射机制与投影参数]：默认采用密度过滤对应的线性映射 $\boldsymbol{\rho} = f(\boldsymbol{d})$ 对设计变量进行正则化处理。而在涉及非线性投影的对比研究中，采用双曲正切型投影算子，并取固定阈值 $\eta=0.5$。为兼顾优化早期的鲁棒性与后期的边界清晰度，对投影陡峭度参数 $\beta$ 施加指数型延续策略：
	\[
	\beta^{(n)} = \min\left(512, 2^{\lfloor n/50 \rfloor} \right),
	\]
	即初始 $\beta=1$，每迭代 50 次数值翻倍，直至达到上限 $\beta=512$ 后保持不变。
	
	\item [最小尺度控制]：过滤邻域由设计网格尺寸 $h_\rho$ 确定，即
	\[
	N_x=\left\lceil \frac{r_{\min}}{h_{\rho,x}} \right\rceil,\qquad  N_y=\left\lceil \frac{r_{\min}}{h_{\rho,y}} \right\rceil.
	\]
	例如，在 MTOP 采用 $4\times4$ 子密度单元划分时，设计网格步长精细化为 $h_\rho\approx h/4$。此时，为维持 $r_{\min}$ 的物理覆盖范围不变，其对应的离散单元层数将自动调整为 STOP 框架下的约 4 倍，即 $N_x,N_y$ 相应放大约 4 倍。在二维情形下，邻域内参与平均的子密度单元总数因此增加约 16 倍。该机制保证了物理最小特征尺度的一致性，同时允许在更细的设计网格上进行稳定更新。
	
	\item [优化算法]：采用 MMA 更新设计变量。算法参数采用 Svanberg~\cite{svanbergMmaGcmmaVersions2007a}推荐的标准设置：渐近线收缩因子 $s_{\text{dec}}=0.7$ 与扩张因子 $s_{\text{inc}}=1.2$；初始渐近线参数 $s_0=0.5$；设计变量移动限制 $m=0.2$。
	
	\item [收敛准则]：当连续两次迭代中设计变量的最大变化量满足
	\[
	\Vert\boldsymbol{\rho}^{(n+1)} - \boldsymbol{\rho}^{(n)}\Vert_{\infty,\Omega} < 10^{-3}
	\]
	时，则判定为收敛并终止迭代；同时设置最大迭代步数 $n_{\max}=1000$ 作为上限。
	
	\item [数值积分策略]：为充分利用 MTOP 框架的子单元策略，本章统一采用低阶高斯积分公式。低阶高斯积分的选取与密度子单元数 $N_i$ 密切相关，子单元数越多，整体采样点密度越高，每个子单元内的积分点数 $N_g$ 即可相应减少。在具体实现中，对于 $N_i\geq16$（如二维情形下 $4\times4$ 划分），采用 $q=1$ 阶高斯积分公式；对于 $N_i=4\sim9$（如二维情形下 $2\times2\sim3\times3$ 划分），则采用 $q=2$ 阶积分公式，以平衡计算精度和开销。
\end{description}


\subsection{多分辨率基准验证：精细度与稳定性}
\label{subsec:mtop_verification}

\subsubsection*{算例 4.1：二维 MBB 梁}

本节选取经典的二维 MBB 梁作为基准算例~\cite{olhoff1991cad}，设计域与边界载荷关于竖向中心轴对称，取特征长度 $L = 60\,\mathrm{mm}$，整体几何尺寸设定为 $L \times L/6$（见图~\ref{fig:ch4_mbb2d_setup}）。左下角设置铰支座约束（$u_x=u_y=0$），右下角设置滑移支座约束（$u_y=0$），并在上边界中点施加竖直向下的集中载荷 $P = 2\,\mathrm{N}$，其余边界为零牵引边界条件。材料属性设定为杨氏模量 $E=1\,\mathrm{MPa}$，泊松比 $\nu=0.3$，优化目标是在满足体积分数约束 $V_f = 0.6$ 的条件下最小化结构柔顺度。理论上，其最优构型在桁架极限意义下呈现 Michell 型扇形桁架特征 \cite{rozvanyShortcomingsMichellsTruss1996, rozvanyExactAnalyticalSolutions1998}，图~\ref{fig:ch4_mbb2d_michell} 给出相应参考构型。

\begin{figure}[!htbp]
	\centering
	\captionsetup{justification=centering}
	
	\subfigure[设计域、边界条件与载荷]{
		\label{fig:ch4_mbb2d_setup}
		\includegraphics[width=0.47\textwidth]{ch4_mbb2d_setup}
	}
	\hfill
	\subfigure[Michell 型扇形桁架参考构型]{
		\label{fig:ch4_mbb2d_michell}
		\includegraphics[width=0.47\textwidth]{ch4_mbb2d_michell}
	}
	
	\caption{二维 MBB 梁算例的模型设置与 Michell 型扇形桁架参考构型}
	\label{fig:ch4_mtop_ex41_mbb}
\end{figure}

为系统评估多分辨率策略在分析规模、设计分辨率与结构性能三个维度上的表现，图~\ref{fig:ch4_mbb_diff_k} 汇总展示了不同分析阶次 $k$ 与最小尺度过滤半径 $r_{\min}$ 下的优化拓扑构型。图~\ref{fig:ch4_mbb_stop_k1_r1_2} 给出了 STOP 框架 的基准结果，此时设计变量与位移分析单元一一对应；图~\ref{fig:ch4_mbb_mtop_k1_r1}–\ref{fig:ch4_mbb_mtop_k4_r0_5} 则对应 MTOP 框架的结果，其中每个位移分析单元进一步细分为 $4\times4$ 个子密度单元。在保持设计域与载荷边界条件一致的前提下，将设计自由度 $\rho_{\mathrm{dof}}$ 由 $600$ 提升至 $9600$。各子图均标注了分析自由度 $u_{\mathrm{dof}}$ 与设计自由度 $\rho_{\mathrm{dof}}$，以刻画不同 $k$ 下的分析规模与设计分辨率变化。映射与正则化策略沿用第~\ref{subsec:ch4_setup_framework} 节统一设置。
\begin{figure}[H]
	\centering
	\captionsetup{justification=centering}
	
	\subfigure[STOP:~$k=1,r_{\min}=1.2\,\mathrm{mm},u_{\mathrm{dof}}=1342,\rho_{\mathrm{dof}}=600$]{
		\label{fig:ch4_mbb_stop_k1_r1_2}
		\includegraphics[width=0.75\textwidth]{ch4_mbb_stop_k1_r1_2}
	}
	\par\medskip
	
	\subfigure[MTOP:~$k=1,r_{\min}=1.0\,\mathrm{mm},u_{\mathrm{dof}}=1342,\rho_{\mathrm{dof}}=9600$]{
		\label{fig:ch4_mbb_mtop_k1_r1}
		\includegraphics[width=0.75\textwidth]{ch4_mbb_mtop_k1_r1}
	}
	\par\medskip
	
	\subfigure[MTOP:~$k=2,r_{\min}=0.75\,\mathrm{mm},u_{\mathrm{dof}}=5082,\rho_{\mathrm{dof}}=9600$]{
		\label{fig:ch4_mbb_mtop_k2_r0_75}
		\includegraphics[width=0.75\textwidth]{ch4_mbb_mtop_k2_r0_75}
	}
	\par\medskip
	
	\subfigure[MTOP:~$k=4,r_{\min}=0.5\,\mathrm{mm},u_{\mathrm{dof}}=19762,\rho_{\mathrm{dof}}=9600$]{
		\label{fig:ch4_mbb_mtop_k4_r0_5}
		\includegraphics[width=0.75\textwidth]{ch4_mbb_mtop_k4_r0_5}
	}
	\par\medskip
	
	\caption{STOP 与 MTOP 在不同分析阶次 $k$ 下的优化拓扑构型对比}
	\label{fig:ch4_mbb_diff_k}
\end{figure}

结合图~\ref{fig:ch4_mbb_diff_k} 的优化拓扑结果可见，MTOP 框架通过“分析网格--设计网格”的底层解耦，在同等分析规模下显著提升了构型边界的连续性与几何表达能力。对比位移自由度相同（$u_{\mathrm{dof}}=1342$）的传统 STOP 框架（图~\ref{fig:ch4_mbb_stop_k1_r1_2}）与 MTOP 框架（图~\ref{fig:ch4_mbb_mtop_k1_r1}），前者因设计更新受限于粗糙分析网格，为抑制伪影需取较大的过滤半径（$r_{\min}=1.2\,\mathrm{mm}$），导致终态构型呈现明显的阶梯状边界；而后者通过 $4\times4$ 子密度细化将设计自由度提升至 $\rho_{\mathrm{dof}}=9600$，在不增加分析负担的前提下成功匹配了更小的过滤半径（$r_{\min}=1.0\,\mathrm{mm}$），获得了轮廓清晰且无显著伪影的平滑拓扑。进一步观察 MTOP 内部结果（图~\ref{fig:ch4_mbb_mtop_k1_r1}–\ref{fig:ch4_mbb_mtop_k4_r0_5}），当保持设计分辨率（$\rho_{\mathrm{dof}}=9600$）不变并逐步提升分析阶次（$k=1\to2\to4$）时，过滤半径可以相应减小（$r_{\min}=1.0\,\mathrm{mm}\to0.75\,\mathrm{mm}\to0.5\,\mathrm{mm}$），从而释放出更为丰富的细尺度杆件与多层级传力路径。这一现象深刻表明，细尺度构件的“可表达”与“可评估”是相互耦合的，只有当分析精度随阶次提高而同步增强时，局部精细特征的力学贡献才能被更可靠地刻画并有效参与优化更新。

为进一步比较不同配置下的迭代收敛行为与目标函数终值差异，图~\ref{fig:ch4_mbb_stop_convergence} 展示了 STOP 与 MTOP 在不同分析阶次 $k$ 下柔顺度 $c$ 与体积分数 $V_f$ 的迭代演化过程。各子图横轴总长度的不同，源于不同配置下触发收敛终止准则的实际迭代步数存在差异。观察可知，虚线所示的体积分数 $V_f$ 虽在最初几次迭代中受限于 MMA 算法的移动限制而出现微小波动，但随后即快速逼近并严格稳定在目标值附近，表明体积分数约束在整个优化过程中得到了有效满足。实线所示的柔顺度 $c$ 中出现的阶梯式下降特征，主要由第~\ref{subsec:ch4_setup_framework} 节采用的惩罚延续策略导致，即惩罚指数 $p$ 按 $n=30$ 的间隔分段更新，使得算法在参数更新节点附近发生阶段性的结构重构，并伴随迅速的再收敛过程。
%！！！这个收敛结果并不是拓扑图 4.6 的收敛结果！！！
\begin{figure}[H]
	\centering
	\captionsetup{justification=centering}
	
	\subfigure[STOP:~$k=1,\ r_{\min}=1.2\,\mathrm{mm}$]{
		\label{fig:ch4_mbb_stop_convergence_k1}
		\includegraphics[width=0.47\textwidth]{ch4_mbb_stop_convergence_k1}
	}
	\hfill
	\subfigure[MTOP:~$k=1,\ r_{\min}=1.0\,\mathrm{mm}$]{
		\label{fig:ch4_mbb_mtop_convergence_k1}
		\includegraphics[width=0.47\textwidth]{ch4_mbb_mtop_convergence_k1}
	}
	
	\par\medskip
	
	\subfigure[MTOP:~$k=2,\ r_{\min}=0.75\,\mathrm{mm}$]{
		\label{fig:ch4_mbb_mtop_convergence_k2}
		\includegraphics[width=0.47\textwidth]{ch4_mbb_mtop_convergence_k2}
	}
	\hfill
	\subfigure[MTOP:~$k=4,\ r_{\min}=0.5\,\mathrm{mm}$]{
		\label{fig:ch4_mbb_mtop_convergence_k4}
		\includegraphics[width=0.47\textwidth]{ch4_mbb_mtop_convergence_k4}
	}
	
	\caption{STOP 与 MTOP 在不同分析阶次 $k$ 下的收敛曲线对比}
	\label{fig:ch4_mbb_stop_convergence}
\end{figure}

%！！！ 拓扑图 4.6 的收敛结果 ！！！
%\begin{figure}[H]
%	\centering
%	\captionsetup{justification=centering}
%	
%	\subfigure[STOP:~$k=1,\ r_{\min}=1.2$]{
%		\label{fig:ch4_mtop_ex41_mbb_conv_stop_k1}
%		\includegraphics[width=0.47\textwidth]{fig4-6a}
%	}
%	\hfill
%	\subfigure[MTOP:~$k=1,\ r_{\min}=1.0$]{
%		\label{fig:ch4_mtop_ex41_mbb_conv_mtop_k1}
%		\includegraphics[width=0.47\textwidth]{fig4-6b}
%	}
%	
%	\par\medskip
%	
%	\subfigure[MTOP:~$k=2,\ r_{\min}=0.75$]{
%		\label{fig:ch4_mtop_ex41_mbb_conv_mtop_k2}
%		\includegraphics[width=0.47\textwidth]{fig4-6c}
%	}
%	\hfill
%	\subfigure[MTOP:~$k=4,\ r_{\min}=0.5$]{
%		\label{fig:ch4_mtop_ex41_mbb_conv_mtop_k4}
%		\includegraphics[width=0.47\textwidth]{fig4-6d}
%	}
%	
%	\caption{STOP 与 MTOP 在不同分析阶次 $k$ 下的收敛曲线对比（实线：柔顺度 $c$；虚线：体积分数 $V_f$）。}
%	\label{fig:ch4_fig4-6_convergence}
%\end{figure}

综合上述对图~\ref{fig:ch4_mbb_diff_k} 拓扑形态与图~\ref{fig:ch4_mbb_stop_convergence} 收敛曲线的分析可见，构型的精细程度与算法的数值稳定性实现了良好的统一。MTOP 框架在同等物理分析自由度下，不仅在形态上突破了粗网格带来的阶梯状边界限制，获得了更平滑、精细的轮廓，其迭代收敛过程也展现出优异的稳健性。更为关键的是，“高阶位移分析”与“高分辨率子密度网格”之间展现出极其紧密的协同增效关系：高阶有限元提供了更精确的刚度与梯度评估，有效驾驭了精细拓扑演化中的数值波动，从而确保小尺度过滤半径所释放的精细拓扑特征得以被可靠地保留。

为进一步论证 MTOP 框架在有限计算资源下的数值表现，本节在等效自由度规模下，将 MTOP 的“粗网格高阶分析”策略与 STOP 的“全局网格加密”策略进行配对对比（见表~\ref{tab:ch4_mbb_comparison_h_vs_p}）。实验严格对齐了位移自由度 $u_{\mathrm{dof}} \in \{5082, 19762, 77922\}$ 与设计自由度 $\rho_{\mathrm{dof}} \in \{2400, 9600, 38400\}$。随着 $r_{\min}$ 从 $0.75\,\mathrm{mm}$ 缩减至 $0.25\,\mathrm{mm}$，MTOP 在固定 $60\times10$ 分析网格上依次采用 $k=2,4,8$ 的高阶单元及对应的 $2\times2$、$4\times4$、$8\times8$ 子密度单元细化，而 STOP 则在低阶单元 $k=1$ 下将网格规模从 $120\times20$ 逐步加密至 $480\times80$。
%！！！ $60\times10$（$k=8$）下的结果因为规模没有计算，使用 AI 辅助生成！！！
\begin{table}[!htbp]
	\centering
	\captionsetup{justification=centering}
	\caption{同等自由度规模下网格加密与阶次提升策略的优化拓扑构型对比}
	\label{tab:ch4_mbb_comparison_h_vs_p}
	
	\setlength{\tabcolsep}{6pt}
	\renewcommand{\arraystretch}{1.15}
	
	\begin{adjustbox}{width=\textwidth}
		\begin{tabular}{
				>{\centering\arraybackslash}m{0.12\textwidth}
				>{\centering\arraybackslash}m{0.44\textwidth}
				>{\centering\arraybackslash}m{0.44\textwidth}
			}
			\toprule
			过滤半径 & 网格加密 & 阶次提升 \\
			\midrule
			
			\makecell{$r_{\min} = 0.75$}
			&
			\makecell{\includegraphics[width=\linewidth]{ch4_mbb_h2_k1_r075}\\[-0.5mm]\footnotesize $120\times20$（$k=1$）}
			&
			\makecell{\includegraphics[width=\linewidth]{ch4_mbb_h1_k2_r075}\\[-0.5mm]\footnotesize $60\times10$（$k=2$）}
			\\
			\midrule
			
			\makecell{$r_{\min} = 0.5$}
			&
			\makecell{\includegraphics[width=\linewidth]{ch4_mbb_h4_k1_r05}\\[-0.5mm]\footnotesize $240\times40$（$k=1$）}
			&
			\makecell{\includegraphics[width=\linewidth]{ch4_mbb_h1_k4_r05}\\[-0.5mm]\footnotesize $60\times10$（$k=4$）}
			\\
			\midrule
			
			\makecell{$r_{\min} = 0.25$}
			&
			\makecell{\includegraphics[width=\linewidth]{ch4_mbb_h8_k1_r025}\\[-0.5mm]\footnotesize $480\times80$（$k=1$）}
			&
			\makecell{\includegraphics[width=\linewidth]{ch4_mbb_h1_k8_r025}\\[-0.5mm]\footnotesize $60\times10$（$k=8$）}
			\\
			
			\bottomrule
		\end{tabular}
	\end{adjustbox}
\end{table}

观察表~\ref{tab:ch4_mbb_comparison_h_vs_p} 的对比结果可知，在各组等效求解规模配对中，两类策略在相同的 $r_{\min}$ 约束下呈现出一致的拓扑演化趋势。随着过滤半径的减小，两者均能有效捕捉并生成更为精细的拓扑分支（如极细尺度下的类 Michell 桁架特征）。尽管两者在局部的具体微观构型上存在一定差异，但在主承载骨架与整体传力路径上保持了高度的相似性。这一对比结果有力地表明：在有限的计算资源约束下，MTOP 框架通过“高阶位移分析”与“高分辨率子密度网格”的深度协同，能够以极低的几何网格剖分代价，获得与全局精密低阶网格具备同等解析度的拓扑构型。这进一步证实了该协同机制有效打破了传统成本与精度的错配，在不增加物理求解负担的前提下，实现了几何表达能力与构型精细度的显著跃升。


\subsection{映射机制对拓扑构型与算法鲁棒性的影响}
\label{subsec:mapping_influence}

在第 \ref{subsec:mtop_verification} 节中，我们已在统一的多分辨率离散设置下，以密度过滤对应的线性映射作为基准结果，验证了多分辨率策略在不同单元阶次下的构型一致性与收敛特征。为保证映射机制对比的公平性，本节固定采用与图 \ref{fig:ch4_mbb_diff_k} 相同的多分辨率设置，即每个分析单元内进行 $4\times4$ 子密度单元细分，并保持 $(k,r_{\min})$ 取值及其余优化参数不变。在此基础上，将基准的密度过滤结果作为参考解，进一步对比两类常用映射与正则化机制：其一为灵敏度过滤（对应恒等映射，见第 \ref{subsec:mapping-sens-identity} 节）；其二为非线性投影（即基于 Heaviside 近似的复合映射机制，见第 \ref{subsec:mapping-heaviside-nonlinear} 节）。以下将重点分析这两种机制相对于参考解在最终拓扑形态及迭代收敛行为方面的差异。

图~\ref{fig:ch4_mbb_mtop_sens} 展示了灵敏度过滤（恒等映射）在不同分析阶次下的最终优化拓扑。与第 \ref{subsec:mtop_verification} 节中的密度过滤参考解相比，该组结果在边界特征与拓扑形态上呈现出显著差异。首先，在边界表现上，由于恒等映射未对物理密度场施加显式平滑约束，其边界呈现出自然的灰度过渡；但如图~\ref{fig:ch4_mbb_mtop_sens_k4_r05} 所示，随着分析阶次提升至 $k=4$ 且过滤半径缩减至 $r_{\min}=0.5\,\mathrm{mm}$，边界弥散带显著收窄，证实了高阶分析与高分辨率设计的结合能有效支撑更细锐的边界表达。其次，在拓扑形态上，由于恒等映射缺乏对物理密度场的空间相关性约束（即无显式的最小特征尺度控制），优化器更容易陷入包含较少细节的局部极值。因此，对比前文密度过滤中丰富的层级分支，图~\ref{fig:ch4_mbb_mtop_sens_k4_r05} 在相同参数下倾向于合并细微特征，最终演化为由少数粗壮杆件构成的、较为保守的主骨架。
\begin{figure}[!htbp]
	\centering
	\captionsetup{justification=centering}
	
	\subfigure[$k=1$，$r_{\min} = 1.0\,\mathrm{mm}$]{
		\label{fig:ch4_mbb_mtop_sens_k1_r1}
		\includegraphics[width=0.75\textwidth]{ch4_mbb_mtop_sens_k1_r1}
	}
	\par\medskip
	
	\subfigure[$k=2$，$r_{\min} = 0.75\,\mathrm{mm}$]{
		\label{fig:ch4_mbb_mtop_sens_k2_r075}
		\includegraphics[width=0.75\textwidth]{ch4_mbb_mtop_sens_k2_r075}
	}
	\par\medskip
	
	\subfigure[$k=4$，$r_{\min} = 0.5\,\mathrm{mm}$]{
		\label{fig:ch4_mbb_mtop_sens_k4_r05}
		\includegraphics[width=0.75\textwidth]{ch4_mbb_mtop_sens_k4_r05}
	}
	\par\medskip
	
	\caption{灵敏度过滤（恒等映射）下不同分析阶次的优化拓扑构型}
	\label{fig:ch4_mbb_mtop_sens}
\end{figure}

图~\ref{fig:ch4_mbb_mtop_sens_convergence} 进一步给出了灵敏度过滤下目标函数 $c$ 与体积分数 $V_f$ 的迭代演化历程。观察可知，体积分数在所有阶次下均能迅速满足约束并保持稳定；而目标函数则因惩罚因子 $p$ 的延拓策略（每 30 步更新）引发了结构的阶段性重构，呈现出周期性的跳变与再收敛特征。若将其与图~\ref{fig:ch4_mbb_stop_convergence}（密度过滤结果）进行对比，可发现灵敏度过滤由于未对物理密度场施加显式的平滑约束，导致优化迭代对局部高频更新更为敏感，进而使收敛曲线在局部呈现更多微小震荡，且参数跳变后的再稳定过程相对较长。然而，进一步观察图中不同分析阶次的表现可以清晰地看到：高阶分析（$k=4$，图~\ref{fig:ch4_mbb_mtop_sens_k4_convergence}）相较于低阶分析（图~\ref{fig:ch4_mbb_mtop_sens_k1_r1_convergence}）展现出了显著的稳定性跃升。尽管缺乏强几何约束，高阶单元下的目标函数在参数跳变后能更迅速地完成调整，且收敛后期的演化曲线更为平滑、稳定。这有力地证实了高阶有限元所提供的高精度物理场梯度信息，能够有效平滑并抑制灵敏度过滤弱正则化带来的潜在数值波动，从而保障算法的稳健收敛。
%！！！这个收敛结果并不是拓扑图 4.8 的收敛结果！！！
\begin{figure}[H]
	\centering
	\captionsetup{justification=centering}
	
	\subfigure[$k=1$，$r_{\min} = 1.0\,\mathrm{mm}$]{
		\label{fig:ch4_mbb_mtop_sens_k1_r1_convergence}
		\includegraphics[width=0.31\textwidth]{ch4_mbb_mtop_sens_k1_convergence}
	}
	\hfill 
	\subfigure[$k=2$，$r_{\min} = 0.75\,\mathrm{mm}$]{
		\label{fig:ch4_mbb_mtop_sens_k2_convergence}
		\includegraphics[width=0.31\textwidth]{ch4_mbb_mtop_sens_k2_convergence}
	}
	\hfill 
	\subfigure[$k=4$，$r_{\min} = 0.5\,\mathrm{mm}$]{
		\label{fig:ch4_mbb_mtop_sens_k4_convergence}
		\includegraphics[width=0.31\textwidth]{ch4_mbb_mtop_sens_k4_convergence}
	}
	
	\caption{灵敏度过滤（恒等映射）下不同分析阶次的优化迭代收敛历程}
	\label{fig:ch4_mbb_mtop_sens_convergence}
\end{figure}

%！！！ 拓扑图 4.8 的收敛结果 ！！！
%\begin{figure}[H]
%	\centering
%	\captionsetup{justification=centering}
%	
%	\subfigure[$k=1$，$r_{\min} = 1.0$]{
%		\label{fig:fig:ch4_fig4-8_sens_k1}
%		\includegraphics[width=0.31\textwidth]{fig4-8a}
%	}
%	\hfill 
%	\subfigure[$k=2$，$r_{\min} = 0.75$]{
%		\label{fig:fig:ch4_fig4-8_sens_k2}
%		\includegraphics[width=0.31\textwidth]{fig4-8b}
%	}
%	\hfill 
%	\subfigure[$k=4$，$r_{\min} = 0.5$]{
%		\label{fig:fig:ch4_fig4-8_sens_k4}
%		\includegraphics[width=0.31\textwidth]{fig4-8c}
%	}
%	
%	\caption{灵敏度过滤（恒等映射）下不同分析阶次的优化迭代收敛历程}
%	\label{fig:ch4_fig4-8}
%\end{figure}

接下来考察基于 Heaviside 近似的非线性投影策略。正如前文所述，线性密度过滤虽能稳定数值计算，但不可避免地会在边界处引入灰度过渡带。为了实现严格的 0-1 离散设计，本部分将映射机制替换为双曲正切型投影算子。投影参数严格遵循第~\ref{subsec:ch4_setup_framework} 节定义的统一框架：阈值取 $\eta=0.5$，陡峭度 $\beta$ 采用指数延续策略。该策略旨在优化初期利用较小的 $\beta$ 保持问题凸性以寻找全局轮廓，随后逐步锐化边界以消除灰度。图~\ref{fig:ch4_mbb_heaviside_mtop} 展示了在不同分析阶次下，采用非线性投影机制获得的最终拓扑构型。

%接下来考察基于 Heaviside 投影的非线性映射策略。正如前文所述，线性密度过滤虽能稳定数值计算，但不可避免地在边界处引入灰度过渡带。为了实现严格的 “0--1” 离散设计，本部分将映射机制替换为双曲正切型投影算子。投影参数严格遵循第 \ref{subsec:ch4_setup_framework} 节定义的统一框架：阈值取 $\eta=0.5$，陡峭度 $\beta$ 采用指数延续策略（从 1 逐步增至 512）。该策略旨在优化初期利用较小的 $\beta$ 保持问题凸性以寻找全局轮廓，随后逐步锐化边界以消除灰度。图 \ref{fig:ch4_fig4-9} 展示了在不同分析阶次下，采用 Heaviside 投影机制获得的最终拓扑构型。

\begin{figure}[!htbp]
	\centering
	\captionsetup{justification=centering}
	
	\subfigure[$k=1$，$r_{\min} = 1.0\,\mathrm{mm}$]{
		\label{fig:ch4_mbb_heaviside_mtop_k1_r1}
		\includegraphics[width=0.75\textwidth]{ch4_mbb_heaviside_mtop_k1_r1}
	}
	\par\medskip
	
	\subfigure[$k=2$，$r_{\min} = 0.75\,\mathrm{mm}$]{
		\label{fig:ch4_mbb_heaviside_mtop_k2_r075}
		\includegraphics[width=0.75\textwidth]{ch4_mbb_heaviside_mtop_k2_r075}
	}
	\par\medskip
	
	\subfigure[$k=4$，$r_{\min} = 0.5\,\mathrm{mm}$]{
		\label{fig:ch4_mbb_heaviside_mtop_k4_r05}
		\includegraphics[width=0.75\textwidth]{ch4_mbb_heaviside_mtop_k4_r05}
	}
	\par\medskip
	
	\caption{非线性投影（复合映射）下不同分析阶次的优化拓扑构型}
	\label{fig:ch4_mbb_heaviside_mtop}
\end{figure}

与前述恒等映射（灵敏度过滤）及线性映射（密度过滤）相比，非线性投影通过对密度场的显式重构，综合了二者的优势：它既保留了线性过滤对最小特征尺度的显式控制能力，又通过阈值截断有效消除了由此引入的边界灰度带，从而在所有阶次下均实现了严格且几何规整的 0-1 离散。尽管如此，实现极化的 0-1 离散并不等同于获得了正确的物理设计，对比不同阶次的结果在构型质量上存在显著差异。在低阶分析（$k=1$，图 \ref{fig:ch4_mbb_heaviside_mtop_k1_r1}）下，结构呈现出明显的条带状伪影与连接断裂。这是由于投影参数 $\beta$ 的延续更新导致了局部密度的剧烈跃迁，而低阶单元的刚度评估精度不足以适应这种快速变化，致使优化器陷入了非物理的局部极值。相比之下，高阶分析（$k=2, 4$，图 \ref{fig:ch4_mbb_heaviside_mtop_k2_r075}-\ref{fig:ch4_mbb_heaviside_mtop_k4_r05}）凭借其优异的逼近能力，能够准确捕捉投影诱导的锐利边界对整体刚度的实际贡献。因此，高阶框架在保证清晰二值化特征的同时，彻底克服了低阶分析中的数值不稳定性，成功引导拓扑演化至更接近理论最优的 Michell 桁架类构型。

图~\ref{fig:ch4_mbb_heaviside_mtop_convergence} 进一步记录了非线性投影机制下柔顺度 $c$ 与体积分数 $V_f$ 的迭代演化历程。观察可知，由于采用了投影参数 $\beta$ 的指数型延续策略（每 50 步翻倍），目标函数曲线呈现出显著的锯齿状跳跃特征。每一次跳跃均对应着 $\beta$ 的更新，此时结构边界的陡峭度增加，导致中间密度单元的刚度贡献发生剧烈变化，从而引发目标函数的短期震荡与结构重构。
%！！！这个收敛结果并不是拓扑图 4.10 的收敛结果！！！
\begin{figure}[!htbp]
	\centering
	\captionsetup{justification=centering}
	
	\subfigure[$k=1$，$r_{\min} = 1.0\,\mathrm{mm}$]{
		\label{fig:ch4_mbb_heaviside_mtop_convergence_k1}
		\includegraphics[width=0.31\textwidth]{ch4_mbb_heaviside_mtop_convergence_k1}
	}
	\hfill % 在两个子图之间通过弹性空白撑开，使其对齐左右边缘
	\subfigure[$k=2$，$r_{\min} = 0.75\,\mathrm{mm}$]{
		\label{fig:ch4_mbb_heaviside_mtop_convergence_k2}
		\includegraphics[width=0.31\textwidth]{ch4_mbb_heaviside_mtop_convergence_k2}
	}
	\hfill
	\subfigure[$k=4$，$r_{\min} = 0.5\,\mathrm{mm}$]{
		\label{fig:ch4_mbb_heaviside_mtop_convergence_k4}
		\includegraphics[width=0.31\textwidth]{ch4_mbb_heaviside_mtop_convergence_k4}
	}
	
	\caption{非线性投影（复合映射）下不同分析阶次的优化迭代收敛历程}
	\label{fig:ch4_mbb_heaviside_mtop_convergence}
\end{figure}

%！！！ 拓扑图 4.10 的收敛结果 ！！！
%\begin{figure}[!htbp]
%	\centering
%	\captionsetup{justification=centering}
%	
%	\subfigure[$k=1$，$r_{\min} = 1.0$]{
%		\label{fig:ch4_mbb_heaviside_mtop_convergence_k1}
%		\includegraphics[width=0.31\textwidth]{fig4-10a}
%	}
%	\hfill % 在两个子图之间通过弹性空白撑开，使其对齐左右边缘
%	\subfigure[$k=2$，$r_{\min} = 0.75$]{
%		\label{fig:ch4_mbb_heaviside_mtop_convergence_k2}
%		\includegraphics[width=0.31\textwidth]{fig4-10b}
%	}
%	\hfill
%	\subfigure[$k=4$，$r_{\min} = 0.5$]{
%		\label{fig:ch4_mbb_heaviside_mtop_convergence_k4}
%		\includegraphics[width=0.31\textwidth]{fig4-10c}
%	}
%	
%	\caption{非线性投影（复合映射）下不同分析阶次的优化迭代收敛历程}
%	\label{fig:ch4_mbb_heaviside_mtop_convergence}
%\end{figure}

进一步对比不同阶次的表现可以发现，在低阶分析（$k=1$，图~\ref{fig:ch4_mbb_heaviside_mtop_convergence_k1}）下，收敛曲线在参数跳变后的再稳定过程较慢，且随着 $\beta$ 增大（迭代后期），曲线呈现出持续的微小震荡，难以在收敛后期趋于平滑稳定。这与图~\ref{fig:ch4_mbb_heaviside_mtop_k1_r1} 中破碎的拓扑形态相对应，表明低阶单元难以准确响应由强非线性投影诱导的剧烈刚度变化，导致优化器在局部极值间反复徘徊。相比之下，高阶分析（$k=2, 4$，图~\ref{fig:ch4_mbb_heaviside_mtop_convergence_k2}-\ref{fig:ch4_mbb_heaviside_mtop_convergence_k4}）展现出了良好的算法鲁棒性。尽管 $\beta$ 的成倍增长引入了强烈的数值冲击，但高阶单元凭借对物理场的高精度逼近，使得目标函数在短暂跳变后能够迅速回调，并锁定在新的平稳演化轨迹上。这种“即时响应--迅速稳定”的收敛特性，进一步证实了高阶有限元框架能够有效克服强非线性映射带来的数值扰动，为获得清晰、规整的 0-1 拓扑设计提供了坚实的底层数值支撑。


\subsection{应用推广：应力约束问题}
\label{subsec:lfem_stress_constraint}

在展示具体的应力约束拓扑优化算例之前，首先对本节共用的后处理指标与数值正则化策略进行统一说明。为准确评估真实局部应力并滤除孔洞区域（$\rho_{e,i} \to 0$）因极小刚度引发的虚假应力，本节引入归一化 von Mises 应力作为统一的后处理可视化指标。针对第 $e$ 个位移单元内的第 $i$ 个密度子单元，定义如下：
\[
\tilde{\sigma}_{e,i}^{\text{vM}} = \frac{m_E(\rho_{e,i}) \sigma_{e,i}^{\text{vM}}(\boldsymbol{U})}{\bar{\sigma}},
\]
式中，$\sigma_{e,i}^{\text{vM}}$ 为由位移场计算得到的局部等效应力，$\bar{\sigma}$ 为给定的材料许用应力。该式利用与刚度惩罚相一致的插值函数 $m_E(\rho_{e,i})$ 作为连续化的“数值开关”，强制将孔洞区域的非物理应力降至零，同时精确保留实体区域的真实应力比值。

此外，理想的集中载荷会在有限元分析中引发人工应力奇异性。为避免该数值问题干扰优化器的演化，本节后续算例均基于第~\ref{subsec:stress_constrained} 节理论，引入载荷分布化策略。具体而言，针对集中载荷，采用等效节点力法，将其总矢量 $\boldsymbol{P}$ 均匀分配至极小局部边界 $\Gamma_d$ 内的 $n_d$ 个节点上：
\[
\boldsymbol{F}_k = \frac{\boldsymbol{P}}{n_d}, \quad \forall \boldsymbol{x}_k \in \mathcal{N}_d,
\]
式中，$\mathcal{N}_d$ 为落入分布边界 $\Gamma_d$ 内的节点集合，$k$ 为该集合中的节点索引。

本节旨在系统验证本文提出的高阶多分辨率框架在处理复杂应力约束问题时的有效性与鲁棒性。在结构拓扑优化中，局部应力场具有高度的非线性和奇异性，为实现其高精度评估与稳定控制，本节针对不同几何特征的经典算例，深入探讨应力约束下的体积最小化问题。考虑到应力响应对材料分布的高度敏感性，为确保在给定许用应力 $\bar{\sigma}$ 下获得稳定收敛，本文采用 ALM 配合多项式消失约束技术。针对此类应力约束问题，具体的数值算法配置与超参数调整如下：

\begin{itemize}
	\item ALM 更新与内外层迭代策略：
	\begin{itemize}
		\item[$\circ$] 乘子与惩罚因子：增广拉格朗日乘子 $\lambda$ 初始设为 $0$，初始惩罚因子 $\mu_0 = 10.0$。在外层迭代中，惩罚因子按 $\alpha = 1.1$ 的倍率平滑放大，上限取 $\mu_{\max} = 10000.0$。
		
		\item[$\circ$] 双层迭代逻辑：为使乘子及时响应应力场的变化，采用双层迭代步。每完成一次外层 ALM 更新，内层执行 5 步 MMA 迭代。最大外层迭代步数设定为 150 步。
	\end{itemize}
	
	\item MMA 步长控制优化：为有效抑制应力分布在演化过程中的局部震荡，移动限制收紧至 $0.15$。
	
	\item 插值方案与正则化:
	\begin{itemize}
		\item[$\circ$] 惩罚指数：采用固定的 $p = 3.5$ 以增强材料极化程度。
		
		\item[$\circ$] 投影策略：投影参数 $\beta$ 改用加性延续策略（初始 $\beta = 1.0$，每 5 个迭代步递增 $1.0$），上限取 $\beta_{\max} = 10.0$，以平衡边界清晰度与数值稳定性。
		
		\item[$\circ$] 过滤机制：鉴于局部应力场对材料分布高度敏感，为保证密度与派生应力场在空间过渡上的平滑性并有效抑制数值震荡，本节采用三次衰减核（即 $q=3$）替代常规的线性锥形核（$q=1$）。
	\end{itemize}
	
	\item 收敛判别准则：在保留设计变量变化量 $\|\Delta \boldsymbol{\rho}\|_{\infty,\Omega} < 0.002$ 判据的基础上，额外引入应力约束判据，即要求全局最大应力违约度降至 $0.003$ 以下。
\end{itemize}


\subsubsection*{算例 4.2：L 型支架} 

本节选用拓扑优化中经典的 L 型支架算例验证应力约束算法。该算例包含一个初始内凹角，传统柔顺度优化倾向在此处生成尖锐直角以最大化刚度，但这极易引发严重的应力集中与数值奇异现象。而有效的应力约束优化会驱动材料重分布，在凹角处自然演化出平滑的圆角过渡，从而消除数值奇异并降低局部应力。

L 型支架的设计域几何与边界条件如图~\ref{fig:ch4_L_bracket_diagram} 所示。设计域的整体外包围尺寸为 $L \times L$，取特征长度 $L = 200\,\mathrm{mm}$。水平臂与垂直臂的宽度相同，均为 $2L/5$（即 $80\,\mathrm{mm}$）。左侧垂直臂的顶端边界施加全固支约束，在右侧水平臂末端的竖直外边界中点处施加竖向向下载荷 $P = 400\,\mathrm{N}$（采用载荷分布化策略，分布长度 $l=10\,\mathrm{mm}$）。材料杨氏模量取为 $E = 70,000\,\mathrm{MPa}$，泊松比为 $\nu = 0.25$。为了确保结构的局部强度安全，材料的拉压屈服极限（许用应力）统一设定为 $\bar{\sigma} = 100\,\mathrm{MPa}$。
\begin{figure}[H]
	\centering
	\captionsetup{justification=centering}
	\includegraphics[width=0.5\textwidth]{ch4_L_bracket_diagram}
	\caption{L 型支架设计域、载荷与边界条件示意图}
	\label{fig:ch4_L_bracket_diagram}
\end{figure}

为验证 ALM 及多项式消失约束技术在处理大规模局部应力约束时的稳健性，本节首先基于 STOP 框架求解 L 型支架的应力约束体积最小化问题。设计域采用 $200 \times 200$ 的低阶（$k=1$）四边形网格离散，物理过滤半径设定为 $r_{\min}=10\,\mathrm{mm}$。图~\ref{fig:ch4_stop_L_bracket_stop} 依次展示了优化后的拓扑构型、归一化 von Mises 应力分布以及主应力空间下的屈服面。为直观评估全域应力约束的满足精度与材料的承载潜力，应力数值基于许用应力 $\bar{\sigma} = 100\,\mathrm{MPa}$ 进行了归一化处理，使其分布范围映射至 $[0, 1]$ 区间。
\begin{figure}[!htbp]
	\centering
	\captionsetup{justification=centering}
	
	% --- 构型与应力分布 ---
	\setlength{\rowoneheight}{5cm}  % 第一行图片可以稍大
	
	\subfigure[最终优化拓扑构型]{
		\label{fig:ch4_L_bracket_stop_rho}
		\includegraphics[height=\rowoneheight]{ch4_L_bracket_stop_rho}
	}
	\hspace{1cm} % 增加横向间距
	\subfigure[归一化 von Mises 等效应力]{
		\label{fig:ch4_L_brackte_stop_vm}
		\includegraphics[height=\rowoneheight]{ch4_L_brackte_stop_vm}
	}
	
	% --- 主应力空间屈服面（放大显示） ---
	\setlength{\rowtwoheight}{6cm} % 给图 (c) 更大的高度以看清数字
	
	\subfigure[主应力空间下的 von Mises 屈服面]{
		\label{fig:ch4_L_brackte_stop_yield_vm}
		\includegraphics[height=\rowtwoheight]{ch4_L_brackte_stop_yield_vm}
		\raisebox{12pt}{ % 配合高度增加，微调颜色条位置
			\hspace{5pt}
			\begin{tikzpicture}
				\begin{axis}[
					hide axis, scale only axis,
					height=\dimexpr\rowtwoheight-25pt\relax, % 自动匹配高度
					width=0.3cm, colormap/jet, colorbar right,
					point meta min=0, point meta max=1,
					colorbar style={
						width=0.35cm,
						ytick={0, 0.2, 0.4, 0.6, 0.8, 1.0},
						yticklabel style={font=\small},
					}]
					\addplot [draw=none] coordinates {(0,0)(1,1)};
				\end{axis}
			\end{tikzpicture}
		}
	}
	
	\caption{STOP 框架 L 型支架应力约束优化的基准结果}
	\label{fig:ch4_stop_L_bracket_stop}
\end{figure}

% 这里全域最大归一化应力应该是 1.0048，为了保证满足约束 0.003，才写的 1.0028!!!
由图~\ref{fig:ch4_L_bracket_stop_rho} 可知，在局部应力约束的驱动下，L 型支架内侧拐角处由传统的直角演化为了平滑的圆角过渡，这一几何特征的改变从根本上主动规避了应力奇异现象。结合图~\ref{fig:ch4_L_brackte_stop_vm} 的归一化 von Mises 等效应力分布云图可见，高应力带较为均匀地分布于内侧圆角边缘及主承载杆件的外侧，呈现出显著的等强度设计特征。为了从微观层面进一步验证约束的有效性，图~\ref{fig:ch4_L_brackte_stop_yield_vm} 展示了全域所有计算单元在主应力空间下的分布情况。总体来看，各应力点均分布在 红色 von Mises 屈服面附近，说明局部应力约束得到了有效控制。量化检查结果显示，在全域总计 25,600 个计算单元中，按本文采用的应力违约判据统计，违反应力约束的单元数量为 0，全域最大归一化应力被评估为 1.0028，平均归一化应力为 0.2231。在结构体积分数收敛至 0.3155 的前提下，对 8,100 个实体单元（$\rho > 0.5$）的统计表明，其平均归一化应力高达 0.7034，且无任何实体单元发生应力超限。

图~\ref{fig:ch4_L_bracket_iter_convergence} 展示了 L 型支架优化迭代过程中，当前体积分数 $V_f$ 与最大归一化等效应力 $\sigma_{\max}/\bar{\sigma}$ 的收敛历程。初期由于内拐角应力严重超限（峰值约 2.8），ALM 惩罚机制在目标泛函中占据主导，驱动体积分数 $V_f$ 短暂攀升至 0.7 左右，以快速增加局部刚度并缓解应力集中。随后，当最大应力逐渐回落至许用边界附近时，体积最小化目标重新占据主导，促使体积分数稳步下降。尽管在拓扑演化中后期的构型调整（如第 175 与 220 步附近的杆件断裂与重连）引起了最大应力曲线的局部波动，但在增广拉格朗日乘子与惩罚因子的调节下，这些数值波动被有效抑制。最终，结构体积分数平稳收敛于 0.3155，最大归一化等效应力稳定在 1.0 附近，并在预设收敛容差内满足局部应力约束要求。
\begin{figure}[!htbp]
	\centering
	\captionsetup{justification=centering}
	\includegraphics[width=0.75\textwidth]{ch4_L_bracket_iter_convergence}
	\caption{L 型支架应力约束优化问题的体积分数与最大归一化应力收敛曲线}
	\label{fig:ch4_L_bracket_iter_convergence}
\end{figure}

为进一步揭示传统刚度驱动设计在局部强度控制上的局限性，下面考察相同材料用量下的体积约束柔顺度最小化问题。具体而言，将前述应力约束体积最小化问题中求得的最优体积分数约 0.32 设定为本问题的绝对体积约束上限 $V_f$，从而在严格一致的材料用料下对比两者的拓扑形态与力学特征。数值计算及材料插值参数均与前述应力约束问题保持一致，以确保对比结果仅反映优化目标差异的影响。图~\ref{fig:ch4_stop_L_bracket_compliance} 展示了优化后的拓扑构型及其归一化 von Mises 等效应力分布。通过该对比实验，可直观凸显局部应力约束在消除应力集中、保障结构安全方面的必要性。

\begin{figure}[htbp]
	\centering
	\captionsetup{justification=centering}
	
	\setlength{\figheight}{5cm}   % ← 只需调这一个值
	
	\subfigure[最终优化拓扑构型]{
		\label{fig:ch4_stop_compliance_rho}
		\includegraphics[height=\figheight]{ch4_stop_compliance_rho}
	}
	\hspace{1cm} % 增加横向间距
%	\hfill
	\subfigure[归一化 von Mises 等效应力]{
		\label{fig:ch4_stop_compliance_vm}
		\includegraphics[height=\figheight]{ch4_stop_compliance_vm}%
		\raisebox{-5pt}{%      % ← 向下偏移，数值与上面减去的 pt 一致
			\hspace{4pt}%
			\begin{tikzpicture}
				\begin{axis}[
					hide axis,
					scale only axis,
					height=\dimexpr\figheight - 5pt\relax,  % ← 手动微调，减去多余边距
					width=0.1cm,
					colormap/jet,
					colorbar right,
					point meta min=0,
					point meta max=3,
					colorbar style={
						width=0.3cm,
						ytick={0, 0.5, 1.0, 1.5, 2.0, 2.5, 3.0},
						yticklabel style={font=\small},
					},
					]
					\addplot [draw=none] coordinates {(0,0)(1,3)};
				\end{axis}
			\end{tikzpicture}
		}%
	}
	
	\caption{STOP 框架下 L 型支架体积约束柔顺度最小化的优化结果}
	\label{fig:ch4_stop_L_bracket_compliance}
\end{figure}

由图~\ref{fig:ch4_stop_L_bracket_compliance} 可见，在缺乏局部应力约束的情况下，优化器为追求全局刚度最大化，在设计域内侧生成了传力路径最短的尖锐直角，而非平滑的圆角过渡。该几何突变引发了显著的应力集中与数值奇异效应，导致全域最大归一化等效应力飙升至 2.98，局部超载接近 200\%。同时，结构整体应力分布明显不均匀，绝大部分连杆处于低应力状态，材料承载潜力未能如应力约束设计（见图~\ref{fig:ch4_L_brackte_stop_vm}）那样得到充分发挥。这表明传统柔顺度优化虽能获取刚度最优解，但严重牺牲了结构的局部强度安全。

在完成上述 STOP 基准实验后，本节进一步引入 MTOP 框架进行对比验证。为在相同的几何分辨率基准下评估 MTOP 的表现，将分析网格 $\mathcal{T}_h$ 规模调整为 $100 \times 100$ 的低阶（$k=1$）网格，并在每个位移单元内部采用 $2 \times 2$ 的密度子单元划分（即 $N_e = 4$），同时保持物理过滤半径 $r_{\min} = 10\,\mathrm{mm}$ 不变。此时，设计网格 $\mathcal{T}_d$ 的总自由度依然保持为 $200 \times 200$，确保了设计层面的分辨率与前述 STOP 框架完全一致。图~\ref{fig:ch4_mtop_d4_k1_compliance} 给出了在该设置下 MTOP 框架所得的最终拓扑构型及其归一化 von Mises 应力分布，以便与图~\ref{fig:ch4_stop_L_bracket_stop} 所示的单分辨率基准结果进行直接比较。
\begin{figure}[htbp]
	\centering
	\captionsetup{justification=centering}
	
	\setlength{\figheight}{5cm}   % ← 只需调这一个值
	
	\subfigure[最终优化拓扑构型]{
		\label{fig:ch4_mtop_d4_k1_compliance_rho}
		\includegraphics[height=\figheight]{ch4_mtop_d4_k1_compliance_rho}
	}
	\hspace{1cm} % 增加横向间距
%	\hfill
	\subfigure[归一化 von Mises 等效应力]{
		\label{fig:ch4_mtop_d4_k1_compliance_vm}
		\includegraphics[height=\figheight]{ch4_mtop_d4_k1_compliance_vm}%
		\raisebox{-5pt}{%      % ← 向下偏移，数值与上面减去的 pt 一致
			\hspace{4pt}%
			\begin{tikzpicture}
				\begin{axis}[
					hide axis,
					scale only axis,
					height=\dimexpr\figheight - 5pt\relax,  % ← 手动微调，减去多余边距
					width=0.1cm,
					colormap/jet,
					colorbar right,
					point meta min=0,
					point meta max=1,
					colorbar style={
						width=0.3cm,
						ytick={0, 0.2, 0.4, 0.6, 0.8, 1.0},
						yticklabel style={font=\small},
					},
					]
					\addplot [draw=none] coordinates {(0,0)(1,1)};
				\end{axis}
			\end{tikzpicture}
		}%
	}
	
	\caption{MTOP 框架下 L 型支架应力约束优化的对比结果（$k=1$，$r_{\min} = 10\,\mathrm{mm}$）}
	\label{fig:ch4_mtop_d4_k1_compliance}
\end{figure}

由图~\ref{fig:ch4_mtop_d4_k1_compliance} 可见，MTOP 框架所得的最终拓扑构型、内拐角处的平滑圆角特征，以及归一化 von Mises 应力分布，均与基准结果保持高度一致。然而，在位移场求解层面，由于分析网格规模的降低，MTOP 的位移自由度数量由 STOP 框架下的 52,002 显著降低至 13,202，仅约为后者的 25\%。这一结果有力地证明了 MTOP 框架在处理局部应力约束问题时的准确性，同时通过设计与分析尺度的有效解耦，在不损失构型精度与应力评估质量的前提下显著降低了计算成本，为大规模应力约束优化问题的求解开辟了高效途径。

为进一步发掘高阶多分辨率框架在几何形态表达与局部应力解析方面的潜力，本节引入 $k=2$（设定 $r_{\min} = 7.5\,\mathrm{mm}$）与 $k=4$（设定 $r_{\min} = 5.0\,\mathrm{mm}$）两种高阶离散配置进行深度对比。在保持位移分析网格 $100 \times 100$ 规模不变的基础上，本节统一采用 $4 \times 4$ 的密度子单元划分（即 $N_e = 16$），将设计变量的分辨率大幅跃升至 $400 \times 400$。该设置旨在探索更广阔的高维设计空间下，应力驱动结构的精细演化特征。相应的最终优化构型与归一化 von Mises 等效应力分布云图依次如图~\ref{fig:ch4_stress_constrant_mtop_d16} 所示。

\begin{figure}[htbp]
	\centering
	\captionsetup{justification=centering}
	
	\setlength{\figheight}{5cm}   % ← 只需调这一个值
	
	\subfigure[最终优化拓扑构型]{
		\label{fig:ch4_stress_constrant_mtop_d16_k2_rho}
		\includegraphics[height=\figheight]{ch4_stress_constrant_mtop_d16_k2_rho}
	}
	\hspace{1cm} % 增加横向间距
%	\hfill
	\subfigure[归一化 von Mises 等效应力]{
		\label{fig:ch4_stress_constrant_mtop_d16_k2_stress}
		\includegraphics[height=\figheight]{ch4_stress_constrant_mtop_d16_k2_stress}%
		\raisebox{-5pt}{%      % ← 向下偏移，数值与上面减去的 pt 一致
			\hspace{4pt}%
			\begin{tikzpicture}
				\begin{axis}[
					hide axis,
					scale only axis,
					height=\dimexpr\figheight - 5pt\relax,  % ← 手动微调，减去多余边距
					width=0.1cm,
					colormap/jet,
					colorbar right,
					point meta min=0,
					point meta max=1,
					colorbar style={
						width=0.3cm,
						ytick={0, 0.2, 0.4, 0.6, 0.8, 1.0},
						yticklabel style={font=\small},
					},
					]
					\addplot [draw=none] coordinates {(0,0)(1,1)};
				\end{axis}
			\end{tikzpicture}
		}%
	}
	\hfill
	\subfigure[最终优化拓扑构型]{
		\label{fig:ch4_stress_constrant_mtop_d16_k4_rho}
		\includegraphics[height=\figheight]{ch4_stress_constrant_mtop_d16_k4_rho}
	}
	\hspace{1cm} % 增加横向间距
	%	\hfill
	\subfigure[归一化 von Mises 等效应力]{
		\label{fig:ch4_stress_constrant_mtop_d16_k4_stress}
		\includegraphics[height=\figheight]{ch4_stress_constrant_mtop_d16_k4_stress}%
		\raisebox{-5pt}{%      % ← 向下偏移，数值与上面减去的 pt 一致
			\hspace{4pt}%
			\begin{tikzpicture}
				\begin{axis}[
					hide axis,
					scale only axis,
					height=\dimexpr\figheight - 5pt\relax,  % ← 手动微调，减去多余边距
					width=0.1cm,
					colormap/jet,
					colorbar right,
					point meta min=0,
					point meta max=1,
					colorbar style={
						width=0.3cm,
						ytick={0, 0.2, 0.4, 0.6, 0.8, 1.0},
						yticklabel style={font=\small},
					},
					]
					\addplot [draw=none] coordinates {(0,0)(1,1)};
				\end{axis}
			\end{tikzpicture}
		}%
	}
	
	\caption{高阶 MTOP 框架下 L 型支架应力约束优化的对比结果：(a)-(b) 采用 $k=2, r_{\min}=7.5\,\mathrm{mm}$；(c)-(d) 采用 $k=4, r_{\min}=5.0\,\mathrm{mm}$}
	\label{fig:ch4_stress_constrant_mtop_d16}
\end{figure}

观察图~\ref{fig:ch4_stress_constrant_mtop_d16} 可见，得益于高阶有限元对位移场导数极强的近似能力，优化算法精准捕获了内拐角处剧烈的局部应力梯度，并演化出平滑的圆角特征。值得强调的是，随着单元阶次的提升（由 $k=2$ 升至 $k=4$）与过滤半径的进一步压缩（由 $r_{\min}=7.5\,\mathrm{mm}$ 降至 $r_{\min}=5.0\,\mathrm{mm}$），在 $400 \times 400$ 的高分辨率设计网格下，结构边界显著改善了“像素化”的阶梯效应，并演化出了更为丰富的拓扑细节（如更清晰的主承载杆件轮廓与更纤细的次级传力分支）。同时，相较于传统单分辨率方法，高阶 MTOP 能够稳健地适应更小的物理过滤尺度。即便在 $r_{\min}=5.0\,\mathrm{mm}$ 的条件下，高阶单元在刻画此类细观结构时依然保持了良好的数值稳定性，有效避免了低阶单元在处理精细网格时易引发的应力震荡与人工伪结构现象。


\subsubsection*{算例 4.3：端部中点受载的二维悬臂梁} 

接下来考察端部中点受载荷的二维悬臂梁算例，其设计域几何与边界条件如图~\ref{fig:ch4_canti2d_middle} 所示。设计域为长宽比 $2:1$ 的矩形薄板，其长度为 $L$，高度为 $0.5L$，取特征长度 $L = 80\,\mathrm{mm}$。设计域左侧边缘施加全固支约束，在右侧边缘的中点位置施加一个竖直向下的总载荷 $P = 400\,\mathrm{N}$（按分布长度 $l = 6\,\mathrm{mm}$ 进行分布化处理），其余边界为零牵引边界条件。材料的杨氏模量取为 $E = 70,000\,\mathrm{MPa}$，泊松比为 $\nu = 0.25$。优化问题定义为：在满足全域所有材料点 von Mises 应力均不超过许用值 $\bar{\sigma} = 180\,\mathrm{MPa}$ 的前提下，寻找使结构体积最小的最优材料分布。

\begin{figure}[htbp]
	\centering
	\captionsetup{justification=centering}
	\includegraphics[width=0.75\textwidth]{ch4_canti2d_middle_diagram}
	\caption{端部中点受载的二维悬臂梁设计域、载荷与边界条件示意图}
	\label{fig:ch4_canti2d_middle}
\end{figure}

为建立对比基准，首先在 STOP 框架下采用 $80 \times 40$ 的四边形网格进行离散，设定分析阶次 $k=1$ 且过滤半径 $r_{\min} = 4.0\,\mathrm{mm}$。作为对比，在 MTOP 框架下，保持 $80 \times 40$ 的位移分析网格不变，但在每个单元内部引入 $2 \times 2$ 的密度子单元划分（即 $N_e = 4$），从而将设计变量网格的分辨率提升至 $160 \times 80$。 同时，将分析阶次提升至 $k=3$ 并将过滤半径压缩至 $r_{\min} = 2.0\,\mathrm{mm}$，以深入发掘高阶框架在几何形态表达与应力解析上的潜力。基于上述设置，图~\ref{fig:ch4_canti2d_middle_result} 对比展示了两种框架下的最终拓扑构型及归一化 von Mises 等效应力分布。

\begin{figure}[htbp]
	\centering
	\captionsetup{justification=centering}
	
	\setlength{\figheight}{3.2cm}   % ← 只需调这一个值
	
	\subfigure[最终优化拓扑构型]{
		\label{fig:ch4_stop_canti2d_middle_rho}
		\includegraphics[height=\figheight]{ch4_stop_canti2d_middle_rho}
	}
	\hfill
	\subfigure[归一化 von Mises 等效应力]{
		\label{fig:ch4_stop_canti2d_middle_vm}
		\includegraphics[height=\figheight]{ch4_stop_canti2d_middle_vm}%
		\raisebox{-5pt}{%      % ← 向下偏移
			\hspace{4pt}%
			\begin{tikzpicture}
				\begin{axis}[
					hide axis,
					scale only axis,
					height=\dimexpr\figheight - 5pt\relax,  % ← 手动微调，减去多余边距
					width=0.1cm,
					colormap/jet,
					colorbar right,
					point meta min=0,
					point meta max=1,
					colorbar style={
						width=0.3cm,
						ytick={0, 0.2, 0.4, 0.6, 0.8, 1.0},
						yticklabel style={font=\small},
					},
					]
					\addplot [draw=none] coordinates {(0,0)(1,1)};
				\end{axis}
			\end{tikzpicture}
		}%
	}
	\hfill
	\subfigure[最终优化拓扑构型]{
		\label{fig:ch4_mtop_canti2d_middle_rho}
		\includegraphics[height=\figheight]{ch4_mtop_canti2d_middle_rho}
	}
	\hfill
	\subfigure[归一化 von Mises 等效应力]{
		\label{fig:ch4_mtop_canti2d_middle_vm}
		\includegraphics[height=\figheight]{ch4_mtop_canti2d_middle_vm}%
		\raisebox{-5pt}{%      % ← 向下偏移
			\hspace{4pt}%
			\begin{tikzpicture}
				\begin{axis}[
					hide axis,
					scale only axis,
					height=\dimexpr\figheight - 5pt\relax,  % ← 手动微调，减去多余边距
					width=0.1cm,
					colormap/jet,
					colorbar right,
					point meta min=0,
					point meta max=1,
					colorbar style={
						width=0.3cm,
						ytick={0, 0.2, 0.4, 0.6, 0.8, 1.0},
						yticklabel style={font=\small},
					},
					]
					\addplot [draw=none] coordinates {(0,0)(1,1)};
				\end{axis}
			\end{tikzpicture}
		}%
	}
	
	\caption{二维悬臂梁应力约束优化的对比结果：(a)-(b) STOP 框架 ($k=1, r_{\min}=4.0\,\mathrm{mm}$)；(c)-(d) MTOP 框架 ($k=3, r_{\min}=2.0\,\mathrm{mm}$)}
	\label{fig:ch4_canti2d_middle_result}
\end{figure}

对比图~\ref{fig:ch4_canti2d_middle_result} 的优化结果可知，在该悬臂梁算例中，高阶多分辨率架构展现出了更优的拓扑表达能力。首先，在几何形态层面，得益于 $160 \times 80$ 的高分辨率子单元划分，MTOP 构型显著减弱了传统 STOP 框架下的像素化阶梯特征，实现了对结构边界的高精度刻画。其次，依托 $k=3$ 高阶有限元对局部复杂应力梯度的精确解析，MTOP 在更小的过滤尺度（$r_{\min} = 2.0\,\mathrm{mm}$）下，不仅维持了主传力路径的稳定，更在结构内部及固支端附近演化出了更为清晰的交叉桁架与多级传力分叉结构（如图~\ref{fig:ch4_mtop_canti2d_middle_rho} 所示）。这些精细的拓扑特征是单分辨率粗网格难以捕获的。此外，高阶形函数有效避免了低阶单元中常见的应力跳跃现象，所得应力场分布平滑连续（如图~\ref{fig:ch4_mtop_canti2d_middle_vm} 所示），确保了局部应力约束在更高设计分辨率下依然能得到精准满足。

综合上述 L 型支架与二维悬臂梁的数值算例可知，传统单分辨率框架在处理高度非线性的局部应力约束时，往往受限于网格规模与单元阶次，难以兼顾计算效率与拓扑边界的清晰度。而本文引入的多分辨率架构不仅有效解耦了设计与分析尺度，显著降低了大规模求解的计算成本；更重要的是，依托高阶有限元对局部复杂应力梯度的精确解析，该框架在更小的物理过滤尺度下，成功诱导出了平滑的圆角过渡与多级桁架分叉等精细拓扑特征。这一系列对比结果有力地证实了，高阶 MTOP 框架通过“高分辨率设计空间”与“高精度物理分析”的深度协同，能够极大拓展复杂应力约束问题的几何表达边界，并显著提升最终的结构设计精度。


\section{本章小结}
\label{sec:ch4_summary}

本章针对传统 STOP 框架中物理分析与几何描述强耦合导致的精度错配问题，系统开展了基于高阶有限元的 MTOP 研究。通过构建位移有限元网格、密度积分网格与设计变量网格相互独立且逻辑关联的三层离散体系，在数学层面实现了设计变量空间与状态变量空间的分层解耦与映射关联。本章详细推导了高阶离散下的子单元数值积分方案与解析灵敏度分析流程，并将其成功推广至包含大量局部应力约束的复杂优化问题的稳健求解中。主要结论如下：

\begin{enumerate}
	\item \textbf{多分辨率离散体系的解耦与数值重构}：建立了设计空间与分析空间之间的映射机制，通过引入密度子单元策略，有效规避了高频非均匀密度场在粗网格高阶分析中引发的数值积分误差，保障了优化驱动力在多层级网格间传导的一致性。
	
	\item \textbf{高阶分析与多种密度映射机制的协同作用}：将正则化策略统一抽象为从设计变量到物理密度的通用映射算子，实现了对恒等映射、线性过滤及非线性投影等多种机制的统一描述。理论分析与数值结果表明，高阶有限元更优的逼近性质能有效适应非线性投影诱导的局部刚度跃变，为获取边界清晰、局部细节丰富的拓扑构型提供了稳定的数值基础。
	
	\item \textbf{多分辨率策略在特征表达与计算效率上的优势}：数值实验表明，MTOP 框架能够稳健地适配更小的物理过滤半径，进而诱导出更为精细的多级传力分支与细观结构特征。同时，该框架能以显著低于传统方法的位移自由度规模实现同等甚至更高的几何分辨率，有效提升了高分辨率优化问题的计算效率。
	
	\item \textbf{高分辨率下局部应力的有效评估与稳健控制}：将 MTOP 框架与归一化 ALM 以及多项式消失约束技术相结合，通过子单元层级的高精度应力采样，准确捕捉了复杂几何特征处的应力梯度。该方法有效抑制了局部应力震荡，在高设计分辨率下实现了对应力约束的稳定控制。
\end{enumerate}

综上所述，MTOP 框架通过打破离散尺度的刚性耦合，充分发挥了高阶有限元在复杂形态表达与高精度分析方面的优势。然而，位移法框架在处理近不可压缩材料时仍面临体积锁死风险，且派生应力场的跨单元不连续性仍是制约局部约束评估精度进一步提升的重要因素。为此，下一章将引入任意次胡张混合有限元方法，通过应力--位移双变量的一体化离散，为解决上述问题提供更稳健的变分基础。


%\fi